\documentclass[acmsmall,screen,nonacm]{acmart}

\usepackage{bcprules}
\usepackage{mathtools}
\usepackage[nameinlink,noabbrev]{cleveref}

\settopmatter{printacmref=false,printccs=false,printfolios=true}
\raggedbottom

\newtheorem{theorem}{Theorem}
\newtheorem{lemma}[theorem]{Lemma}

\newcounter{RuleRef}
\newcommand{\ruledef}[1]{%
  \refstepcounter{RuleRef}\textsc{#1}\label{rule:#1}}
\newcommand{\ruleuse}[1]{%
  \hyperref[rule:#1]{\textsc{#1}}}
\newcommand{\ruleref}[1]{%
  \textsc{(\hyperref[rule:#1]{#1})}}

\newcommand{\lambdap}{\ensuremath{\lambda_{\mathsf p}}}
\newcommand{\ctxempty}{\varnothing}
\newcommand{\sing}[1]{\{#1\}}
\newcommand{\pairty}[4]{\langle #1\!:\!#2;\,#3\!:\!#4\rangle}
\newcommand{\interty}[2]{#1\mathbin{\wedge}#2}
\newcommand{\unionty}[2]{#1\mathbin{\vee}#2}
\newcommand{\mergeop}{\mathbin{\bowtie}}
\newcommand{\joinop}{\mathbin{\sqcup}}
\newcommand{\pairval}[3]{\langle #1;\,#2=#3\rangle}
\newcommand{\subst}[3]{#1[#2/#3]}
\newcommand{\wf}{\mathrel{\mathsf{wf}}}
\newcommand{\final}{\mathsf{final}}
\newcommand{\steps}{\longrightarrow^{*}}
\newcommand{\sstep}{\longrightarrow}
\newcommand{\resolve}{\Downarrow}
\newcommand{\loc}[1]{\mathsf{loc}(#1)}
\newcommand{\typeref}[1]{\mathsf{type}(#1)}
\newcommand{\referentof}[1]{\mathsf{ref}(#1)}
\newcommand{\realizes}{\Vdash}
\newcommand{\coerce}{\mathrel{\rightsquigarrow}}
\newcommand{\bodycl}[5]{%
  \mathsf{Body}_{#1}(#2\!:\!#3\vdash #4:#5)}
\newcommand{\membercl}[5]{%
  \mathsf{Member}_{#1}(#2\!:\!#3\vdash #4<:#5)}
\newcommand{\resolvepi}[2]{\mathsf{resolve}_{#1}(#2)}
\newcommand{\compilepi}[2]{\llbracket #2\rrbracket_{#1}}
\newcommand{\action}[2]{#1(#2)}
\newcommand{\forcemem}[2]{\mathsf{inst}_{\mathsf{mem}}(#1,#2)}
\newcommand{\forcecod}[2]{\mathsf{inst}_{\mathsf{cod}}(#1,#2)}
\newcommand{\closure}[3]{\langle #1,#2\rangle}
\newcommand{\judgheader}[3][]{%
  \par\addvspace{0.75em}%
  \noindent\textbf{#2}%
  \if\relax\detokenize{#1}\relax\else\quad #1\fi
  \hfill{\setlength{\fboxsep}{4pt}\setlength{\fboxrule}{0.5pt}%
    \fbox{\(\strut #3\)}}\par
  \nobreak\addvspace{0.85em}%
}
\newcommand{\plateheader}[2]{%
  \par\addvspace{1.1em}%
  \noindent\textbf{#1}%
  \if\relax\detokenize{#2}\relax\else
    \hfill{\setlength{\fboxsep}{3pt}\setlength{\fboxrule}{0.5pt}%
      \fbox{\(\strut #2\)}}%
  \fi\par
  \nobreak\addvspace{0.65em}%
}
\newcommand{\rulesetup}{%
  \small
  \smallrulenames
  \setlength{\afterruleskip}{2em plus 0.2em minus 0.1em}%
  \setlength{\labelminsep}{0.6em}%
  \renewcommand{\arraystretch}{1.15}%
  \renewcommand{\andalso}{\qquad}%
}

\title{Semantic Coercions for Paths through Dependent Pairs}
\author{Oliver Bra\v{c}evac}
\author{Martin Odersky}
\renewcommand{\shortauthors}{Bra\v{c}evac and Odersky}
\authorsaddresses{}

\begin{document}

\begin{abstract}
We study a small call-by-value calculus of path-dependent types based on
immutable dependent pairs. A pair contains a first-component location and
either a term member with a proper-type signature or a type member with an
interval signature; the member signature may mention the first component.
Proper types also include binary intersections and unions with the standard
meet and join rules. A proof-relevant merge judgment can retain two proper
types as an opaque intersection or recursively normalize aligned pair views.
The pair case requires the same outer label, recursively merges the first
components, and merges the member signatures under the resulting
first-component type. Interval members preserve an equal lower bound or join
distinct lowers, while recursively merging their upper bounds. One subtyping
rule turns a chosen merge plan into a subtype from the original intersection.
An intersection records two semantic realizations of one location; a union
records a tagged realization of one arm. Neither adds a value form, cast, or
machine transition.
The central subtyping rule is covariant in both pair components and compares
the member signatures under the source first-component type, so its dependent
premise cannot be interpreted until a concrete pair realization supplies the
location bound by that premise.

The soundness proof relates paths to their run-time referents, which are
locations or stored type definitions. A realization derivation records the
signature supported by a referent. Each subtyping derivation is interpreted
as a store-indexed coercion that maps a realization of its source signature
to a realization of its target signature at the same referent. Dependent
premises are retained with their closing substitutions until the corresponding
location is known. In the pair-covariance case, inversion supplies the stored first
component; the member premise is interpreted at that location, and the
resulting coercion is applied to the member referent. Coercion application is
well founded lexicographically by referent rank and coercion size; the separate
merge action is structurally recursive on its finite plan. We give the complete
calculus and operational semantics, the coercion interpretation, the
machine-typing invariant, and mechanized proofs of progress and preservation.
\end{abstract}

\keywords{path-dependent types, dependent pairs, intersection types, union
types, subtyping, semantic coercions, type soundness, Lean}

\maketitle

\section{Introduction}

Path-dependent types associate type information with stable term paths. DOT
calculi place this mechanism in a rich object calculus with recursive self
types, refinements, intersections, and abstract type members
\citep{amin2014foundations,rompf2016dot}. The interaction among those
features has led to several distinct soundness arguments. The calculus in
this paper, \lambdap{}, isolates a smaller question: how much of paths and
abstract type selection can be based on ordinary dependent pairs?

A \lambdap{} value is an abstraction or an immutable pair
\(\pairval{y}{a}{\delta}\). Its first component is a store location. The
second component is either another location, representing a term member, or
a type definition, representing a type member. Its signature may refer to the
first component. Paths consist of variables, first projections, and labelled
selections. When the outer pair has a different label, selection continues
through its first component, so nested pairs can represent record extension.
Singleton types \(\sing{p}\) record path identity. A type selection \(p.A\)
has the lower and upper bounds declared by the corresponding pair signature;
the run-time pair itself stores a concrete type definition. The language also
contains dependent functions, declarative subtyping with an explicit
transitivity rule, binary meets and joins, and bound-consistency premises for
type intervals. Unions are proper types, not term-level sums: the language has
no injection or case term.

Write \(\Gamma\vdash\tau<:\tau'\) when signature \(\tau\) is a subtype of
\(\tau'\) under typing context \(\Gamma\). The difficult rule is covariance
of dependent pairs. Here \(S,S'\) range over proper types, while
\(\tau,\tau'\) range over member signatures, which may be either proper
types or type intervals. The type \(\pairty{x}{S}{a}{\tau}\) binds \(x\) in
\(\tau\), where \(x\) denotes the first component and \(a\) is the member
label:
\begingroup
\rulesetup
\typicallabel{s-pair}
\infrule[\ruleuse{s-pair}]
  {\Gamma\vdash S<:S'
   \andalso
   \Gamma,x:S\vdash \tau<:\tau'}
  {\Gamma\vdash
    \pairty{x}{S}{a}{\tau}<:\pairty{x}{S'}{a}{\tau'}}
\endgroup
The source pair supplies a component of type \(S\), so \(S\) is the natural
assumption for comparing the members. At run time, however, the member
signature is opened with the concrete location \(y\) stored in the pair. A
proof of preservation must use the second premise at that location and must
relate source paths to the referents found by run-time resolution.

We interpret subtyping derivations using store-indexed coercions. Fix a store
\(\sigma\) and a closing substitution \(\rho\) that maps the variables of
\(\Gamma\) to locations in \(\sigma\). A path resolves to a \emph{referent},
which is either a location or a stored type definition. We write
\(\sigma\realizes r:\tau\) when referent \(r\) realizes signature \(\tau\).
A derivation of \(\Gamma\vdash\tau<:\tau'\) is interpreted as a coercion
\[
  C:\sigma\vdash\rho(\tau)\coerce\rho(\tau').
\]
Applying \(C\) to a derivation of
\(\sigma\realizes r:\rho(\tau)\) produces a derivation of
\(\sigma\realizes r:\rho(\tau')\). Transitivity is interpreted by coercion
composition. The singleton and selection cases retain resolution derivations
and lower- or upper-bound coercions needed when the coercion is applied.
Dependent
premises under function and pair binders are retained until a concrete binder
location is available.

For \ruleref{s-pair}, interpretation produces a coercion for \(S<:S'\) and
suspends the derivation of \(\tau<:\tau'\) with its closing substitution.
Applying the pair-covariance coercion at a location instantiates the suspended premise
at the stored first component \(y\) and applies the resulting member coercion
to the member referent (\cref{sec:coercion-pair}). Although that coercion is
not structurally contained in the pair-covariance coercion, its application decreases
the allocation rank, because both referents stored in a pair precede the
pair cell in the immutable store (\cref{sec:coercion-termination}).

For \ruleref{s-merge}, interpretation retains the selected finite merge plan
together with \(\mathcal R\). Applying the resulting coercion to an
intersection realization invokes a separate structural action on that plan.
At a pair node the action first merges the concrete first-component evidence,
extends \(\mathcal R\) at that location, and interprets the member subplan in
the extended environment. This is how a dependent later member remains tied
to the record prefix produced by recursive normalization.

The coercions are metatheoretic proof objects. Programs and
machine states contain the source terms given in \cref{sec:calculus};
execution performs no casts or type checks. Typed intermediate calculi such
as System FC instead place equality evidence and casts in the intermediate
program and erase them before execution \citep{sulzmann2007systemfc}.

The contributions are:
\begin{itemize}
\item a small calculus combining arbitrary stable paths, singleton types,
  binary intersections and unions, abstract type selections, dependent
  functions, and dependent pairs with term and type members, including a
  recursive, binder-aware judgment for normalizing aligned record spines;
\item a store-indexed coercion interpretation of its full declarative
  subtyping relation, including the explicit transitivity rule and covariance
  of both components of dependent pairs, together with a structural semantic
  action for merge plans;
\item a well-founded definition of coercion application based on the allocation order
  of immutable stores; and
\item a Lean 4 mechanization of typed path resolution, subtyping
  interpretation, coercion application, machine typing, progress,
  preservation, and
  non-stuckness of every finite execution prefix.
\end{itemize}

The result is type safety for the calculus as presented. It does not establish
normalization, decidability of subtyping, or admissibility of transitivity.

\section{The calculus}
\label{sec:calculus}

\lambdap{} isolates path-dependent types in a call-by-value calculus with
immutable dependent pairs. Terms sequence computations with
\(\mathbf{let}\); applications take paths as operands, and pair components
are variables or type definitions. The operational semantics combines path
resolution with a CK machine.

\subsection{Source language}
\label{sec:syntax}

Proper types and member signatures are defined mutually. In the grammar,
\(S,T,U\) range over proper types; \(\tau,\tau'\) range over proper
or interval signatures; \(p,q\) range over paths; \(t,u\) range over
terms; \(v\) ranges over values; \(x,y,z\) range over variables; \(a,b\)
range over labels; \(\delta\) ranges over member definitions; and \(\Gamma\)
ranges over typing contexts. Primes and numeric subscripts distinguish
metavariables of the same category. Uppercase labels such as \(A\) are a
typographical convention for type members and do not form a separate label
sort.

A member signature is either a proper type or an interval. A term-member
definition is a variable that names a location at run time and therefore has
a proper-type signature. A type-member definition is a type and therefore
has an interval signature.
Labels occupy one syntactic name space.

The binary proper types \(\interty{S}{T}\) and \(\unionty{S}{T}\) carry the
standard three meet and three join rules. A meet is introduced from two views
of one source and eliminates to either component; either component injects
into a join, and a join eliminates when both arms have the same target. The
joins are proof-oriented proper types rather than term-level sums: there are
no injection values or case terms.

Aligned record normalization is described by a separate proof-relevant
judgment
\[
  \Gamma\vdash \tau_1\mergeop\tau_2\Longrightarrow\tau.
\]
Equal signatures may be kept unchanged, while any two proper types may be
left as their opaque intersection. The recursive pair plan is
\[
  \frac{
    \Gamma\vdash S_1\mergeop S_2\Longrightarrow S
    \qquad
    \Gamma,x:S\vdash\tau_1\mergeop\tau_2\Longrightarrow\tau
  }{
    \Gamma\vdash
    \pairty{x}{S_1}{a}{\tau_1}\mergeop
    \pairty{x}{S_2}{a}{\tau_2}
    \Longrightarrow
    \pairty{x}{S}{a}{\tau}
  }.
\]
The member plan is deliberately checked under the \emph{merged}
first-component type \(S\). Consequently later members may depend on the
record prefix that the first premise has just normalized. Repeated use of
this rule traverses an arbitrarily long aligned record spine, and it can
merge both the prefix and the outer member in one derivation.

For interval members a small lower-join judgment either preserves a common
lower bound or forms a union:
\[
  \Gamma\vdash L\joinop L\Longrightarrow L,
  \qquad
  \Gamma\vdash L_1\joinop L_2\Longrightarrow\unionty{L_1}{L_2}.
\]
Together with an upper merge it yields
\[
  \frac{
    \Gamma\vdash L_1\joinop L_2\Longrightarrow L
    \qquad
    \Gamma\vdash U_1\mergeop U_2\Longrightarrow U
  }{
    \Gamma\vdash (L_1..U_1)\mergeop(L_2..U_2)
      \Longrightarrow L..U}.
\]
Finally, one source subtyping rule executes a chosen plan:
\[
  \frac{\Gamma\vdash S\mergeop T\Longrightarrow U}
       {\Gamma\vdash\interty{S}{T}<:U}.
\]
This judgment is plan directed rather than deterministic. For aligned pair
types one may stop at the opaque-intersection plan or choose the recursive
pair plan; the calculus does not claim a unique normal form.

The four earlier one-cell merge principles are now derived instances. They
respectively choose: an intersecting first plan and an unchanged member; an
unchanged first plan and an intersecting proper member; an unchanged first
plan with a shared interval lower and intersected upper; or an unchanged
first plan with a union lower and intersected upper. Thus they still derive
\[
\begin{gathered}
  \interty{\pairty{x}{S}{a}{\tau}}{\pairty{x}{T}{a}{\tau}}
    <:\pairty{x}{\interty{S}{T}}{a}{\tau},
  \\
  \interty{\pairty{x}{S}{a}{T}}{\pairty{x}{S}{a}{U}}
    <:\pairty{x}{S}{a}{\interty{T}{U}},
  \\
  \interty{\pairty{x}{S}{A}{L..U}}{\pairty{x}{S}{A}{L..V}}
    <:\pairty{x}{S}{A}{L..\interty{U}{V}},
  \\
  \interty{\pairty{x}{S}{A}{L_1..U_1}}
           {\pairty{x}{S}{A}{L_2..U_2}}
    <:\pairty{x}{S}{A}
       {\unionty{L_1}{L_2}..\interty{U_1}{U_2}}.
\end{gathered}
\]
The shared-lower plan remains useful because it preserves \(L\) exactly,
instead of producing \(\unionty{L}{L}\).

A merge plan proves subtyping, not well-formedness. To use its target as a
term type through \ruleref{t-sub}, one must separately prove that target well
formed. In particular, an interval target may require
\(\unionty{L_1}{L_2}<:\interty{U_1}{U_2}\), including the cross-bound
comparisons \(L_1<:U_2\) and \(L_2<:U_1\).
There is no dedicated term-level intersection or union introduction: existing
subsumption is the only way to assign either type to a term. In particular,
path synthesis is not extended with rules that project through intersections
or unions. Recursive pair plans require identical labels, member sorts, and
record order at every pair node they traverse. They do not merge unequal
labels, reorder fields, or provide unrestricted distributivity.

The binder in a function scopes over its codomain. The binder \(x\) in
\(\pairty{x}{S}{a}{\tau}\) scopes over the member signature \(\tau\); it
denotes the first component of the pair. We write
\(\subst{\tau}{p}{x}\) for capture-avoiding opening. Singleton types
\(\sing{p}\) record term identity. An abstract type selection \(p.A\) denotes an
unknown type constrained by the bounds declared for \(A\).

The paper uses names, with \(\Gamma(x)=T\) for context lookup, and treats
well-scopedness of the syntax in each judgment as implicit. Substitution and
renaming are capture avoiding; beneath a binder they act on the free variables
and leave the bound variable fixed.

\subsection{Static semantics}
\label{sec:static-semantics}

The static semantics consists of path synthesis
\(\Gamma\vdash p\Rightarrow\tau\), signature subtyping
\(\Gamma\vdash\tau<:\tau'\), signature well-formedness
\(\Gamma\vdash\tau\ \wf\), and term typing \(\Gamma\vdash t:T\). Path
synthesis produces either a proper type or an interval and admits no
subsumption. Term typing assigns a singleton type to a path and provides
subsumption separately. The semantic interpretation is derivation directed:
its result may depend on the chosen derivation of a judgment.

The matching selection rule opens the member at the receiver's first
projection. A pair can extend a chain of earlier pairs: if its label differs
from the requested label, \ruleref{p-miss} continues at \(p.1.a\). Run-time
path resolution uses the same convention. Function domains are contravariant,
and codomains are compared under the target domain. Dependent pairs are
covariant in their first component and their member; the member comparison
is made under the source first-component type and applies uniformly to
proper-type and interval signatures.

Binary intersections use the three meet rules: a source below both components
is below their intersection, and an intersection is below either component.
Binary unions dually use two injections and an elimination rule whose two
premises share a target.
Rule \ruleref{s-merge} turns a proof-relevant merge plan into subtyping from
an intersection. The four plan rules either preserve an identical signature,
form an opaque proper-type intersection, recursively merge same-label pair
views, or merge interval signatures. The pair rule checks its member plan
under the recursively merged first-component type. The two lower-join plans
preserve an identical interval lower or form its union. No merge plan changes
the precise type synthesized by path typing. As with subtyping generally,
\ruleref{s-merge} does not assert that its target is well formed. Any use of
that target in term subsumption supplies a separate well-formedness
derivation, including \ruleref{wf-bounds} for every merged interval.

Well-formed intervals have consistent bounds. Term paths receive singleton
types. A selected type is well formed whenever path synthesis resolves its
member to an interval, including lookup through earlier pair components.
Pair introduction exposes exact information: a term member
receives the singleton type of its location, and a type definition receives
the exact interval \(T..T\). Subtyping subsequently hides either component. In
\ruleref{t-let}, the result type \(T\) is formed in the outer context and does
not depend on \(x\).

\subsection{Runtime configurations and transitions}
\label{sec:machine}

A store is an immutable sequence of values. Source variables and allocated
locations share the metavariables \(x,y,z\); their role is determined by the
judgment in which they occur. A value appended at a fresh location may refer
only to older cells. We write \(\sigma(x)=v\) for lookup,
\(\sigma\mathbin{\cdot}(z\mapsto v)\) for extension at a fresh location
\(z\), and
\(\mathsf{dom}(\sigma)=\{i\mid i<|\sigma|\}\) for the allocated locations.
Machine configurations are well scoped by their store, so every free runtime
variable names an allocated location. Runtime terms therefore use the source
grammar.

Intersections and unions are absent from values, referents, continuations, and
machine states except as types decorating the metatheory. They add no
path-resolution or transition rule, and execution constructs neither
intersection packages nor union injections and performs no casts. Tagged
union evidence appears only in the semantic realization relation. Thus the
runtime syntax and operational semantics below are unchanged by these proper
type connectives.

Path resolution yields either a location or a stored type; machine rules that
require a location stipulate a result of the form \(\loc{x}\).

The CK machine evaluates a let's bound computation before its body. A location
returns through a suspended let body without allocation. A syntactic value is
allocated and the body is opened at its fresh location.

\clearpage
\begingroup
\setlength{\abovedisplayskip}{5pt plus 1pt minus 1pt}
\setlength{\belowdisplayskip}{5pt plus 1pt minus 1pt}
\setlength{\intextsep}{4pt}
\captionsetup{skip=6pt}
\rulesetup
\setlength{\afterruleskip}{0.3em plus 0.06em minus 0.03em}

\begin{figure}[H]
\noindent
\begin{minipage}[t]{\linewidth}
\plateheader{Source syntax}{}
\noindent
\begin{minipage}[t]{0.48\linewidth}
\begin{tabular*}{\linewidth}{@{}r@{\ }l@{\extracolsep{\fill}}r@{}}
\(x,y,z\) && \textbf{Term variable}
\\
\(a,b,A\) && \textbf{Member label}
\\[0.35em]
\(p,q\coloneqq\) && \textbf{Path}
\\
  & \(x\) & variable
\\
  & \(p.1\) & first projection
\\
  & \(p.a\) & member selection
\\[0.35em]
\(S,T,U\coloneqq\) && \textbf{Proper type}
\\
  & \(\top\) & top
\\
  & \(\bot\) & bottom
\\
  & \((x:S)\to T\) & dependent function
\\
  & \(\pairty{x}{S}{a}{\tau}\) & dependent pair
\\
  & \(\interty{S}{T}\) & binary intersection
\\
  & \(\unionty{S}{T}\) & binary union
\\
  & \(\sing{p}\) & singleton
\\
  & \(p.A\) & type selection
\\[0.35em]
\(\tau,\tau'\coloneqq\) && \textbf{Member signature}
\\
  & \(T\) & proper type
\\
  & \(S..T\) & interval
\end{tabular*}
\end{minipage}\hfill
\begin{minipage}[t]{0.48\linewidth}
\begin{tabular*}{\linewidth}{@{}r@{\ }l@{\extracolsep{\fill}}r@{}}
\(\delta\coloneqq\) && \textbf{Member definition}
\\
  & \(z\) & term member
\\
  & \(T\) & type member
\\[0.35em]
\(t,u\coloneqq\) && \textbf{Term}
\\
  & \(p\) & path
\\
  & \(\lambda(x:S).t\) & abstraction
\\
  & \(\pairval{y}{a}{\delta}\) & pair
\\
  & \(p\,q\) & application
\\
  & \(\mathbf{let}\ x=t\ \mathbf{in}\ u\) & let
\\[0.35em]
\(v\coloneqq\) && \textbf{Value}
\\
  & \(\lambda(x:S).t\) & abstraction
\\
  & \(\pairval{y}{a}{\delta}\) & pair
\\[0.35em]
\(\Gamma\coloneqq\) && \textbf{Typing context}
\\
  & \(\ctxempty\) & empty
\\
  & \(\Gamma,x:T\) & term binding
\end{tabular*}
\end{minipage}
\end{minipage}\par
\begin{minipage}[t]{\linewidth}
\plateheader{Path synthesis}{\Gamma\vdash p\Rightarrow\tau}
\typicallabel{p-miss}
\noindent
\begin{minipage}[t]{0.47\linewidth}
\infrule[\ruledef{p-var}]
  {\Gamma(x)=T}
  {\Gamma\vdash x\Rightarrow T}

\infrule[\ruledef{p-fst}]
  {\Gamma\vdash p\Rightarrow\pairty{x}{S}{a}{\tau}}
  {\Gamma\vdash p.1\Rightarrow S}
\end{minipage}\hfill
\begin{minipage}[t]{0.51\linewidth}
\infrule[\ruledef{p-sel}]
  {\Gamma\vdash p\Rightarrow\pairty{x}{S}{a}{\tau}}
  {\Gamma\vdash p.a\Rightarrow\subst{\tau}{p.1}{x}}

\infrule[\ruledef{p-miss}]
  {\Gamma\vdash p\Rightarrow\pairty{x}{S}{b}{\tau'}
   \\
   \Gamma\vdash p.1.a\Rightarrow\tau
   \andalso
   a\ne b}
  {\Gamma\vdash p.a\Rightarrow\tau}
\end{minipage}
\end{minipage}

\plateheader{Signature well-formedness}{\Gamma\vdash\tau\ \wf}
\noindent
\begin{minipage}[t]{0.47\linewidth}
\typicallabel{wf-single}
\infax[\ruledef{wf-bot}]
  {\Gamma\vdash\bot\ \wf}

\infax[\ruledef{wf-top}]
  {\Gamma\vdash\top\ \wf}

\infrule[\ruledef{wf-single}]
  {\Gamma\vdash p\Rightarrow T}
  {\Gamma\vdash\sing{p}\ \wf}

\infrule[\ruledef{wf-select}]
  {\Gamma\vdash p.A\Rightarrow T..U
   \andalso
   \Gamma\vdash T<:U}
  {\Gamma\vdash p.A\ \wf}

\infrule[\ruledef{wf-inter}]
  {\Gamma\vdash S\ \wf
   \andalso
   \Gamma\vdash T\ \wf}
  {\Gamma\vdash\interty{S}{T}\ \wf}

\infrule[\ruledef{wf-union}]
  {\Gamma\vdash S\ \wf
   \andalso
   \Gamma\vdash T\ \wf}
  {\Gamma\vdash\unionty{S}{T}\ \wf}
\end{minipage}\hfill
\begin{minipage}[t]{0.51\linewidth}
\typicallabel{wf-bounds}
\infrule[\ruledef{wf-fun}]
  {\Gamma\vdash S\ \wf
   \andalso
   \Gamma,x:S\vdash T\ \wf}
  {\Gamma\vdash (x:S)\to T\ \wf}

\infrule[\ruledef{wf-pair}]
  {\Gamma\vdash S\ \wf
   \andalso
   \Gamma,x:S\vdash\tau\ \wf}
  {\Gamma\vdash\pairty{x}{S}{a}{\tau}\ \wf}

\infrule[\ruledef{wf-bounds}]
  {\Gamma\vdash S\ \wf
   \andalso
   \Gamma\vdash T\ \wf
   \\
   \Gamma\vdash S<:T}
  {\Gamma\vdash S..T\ \wf}
\end{minipage}

\caption{Source syntax and static semantics: source syntax, path
synthesis, and signature well-formedness.}
\Description{The complete source grammar, all four path-synthesis
rules, and all nine signature well-formedness rules.}
\label{fig:source-syntax}
\label{fig:path-typing}
\label{fig:well-formedness}
\end{figure}

\clearpage
\begin{figure}[H]
\ContinuedFloat
\plateheader{Lower-bound join plan}
  {\Gamma\vdash L_1\joinop L_2\Longrightarrow L}
\noindent
\begin{minipage}[t]{0.47\linewidth}
\infax[\ruledef{j-same}]
  {\Gamma\vdash L\joinop L\Longrightarrow L}
\end{minipage}\hfill
\begin{minipage}[t]{0.51\linewidth}
\infax[\ruledef{j-union}]
  {\Gamma\vdash L_1\joinop L_2
    \Longrightarrow\unionty{L_1}{L_2}}
\end{minipage}

\plateheader{Aligned signature merge plan}
  {\Gamma\vdash\tau_1\mergeop\tau_2\Longrightarrow\tau}
\noindent
\begin{minipage}[t]{0.47\linewidth}
\infax[\ruledef{m-same}]
  {\Gamma\vdash\tau\mergeop\tau\Longrightarrow\tau}

\infax[\ruledef{m-inter}]
  {\Gamma\vdash S\mergeop T\Longrightarrow\interty{S}{T}}
\end{minipage}\hfill
\begin{minipage}[t]{0.51\linewidth}
\infrule[\ruledef{m-intv}]
  {\Gamma\vdash L_1\joinop L_2\Longrightarrow L
   \andalso
   \Gamma\vdash U_1\mergeop U_2\Longrightarrow U}
  {\Gamma\vdash
    (L_1..U_1)\mergeop(L_2..U_2)\Longrightarrow L..U}
\end{minipage}

\infrule[\ruledef{m-pair}]
  {\Gamma\vdash S_1\mergeop S_2\Longrightarrow S
   \andalso
   \Gamma,x:S\vdash\tau_1\mergeop\tau_2\Longrightarrow\tau}
  {\Gamma\vdash
   \pairty{x}{S_1}{a}{\tau_1}\mergeop
   \pairty{x}{S_2}{a}{\tau_2}
   \Longrightarrow
   \pairty{x}{S}{a}{\tau}}

\plateheader{Signature subtyping}{\Gamma\vdash\tau<:\tau'}
\noindent
\begin{minipage}[t]{0.47\linewidth}
\typicallabel{s-trans}
\infax[\ruledef{s-refl}]
  {\Gamma\vdash\tau<:\tau}

\infrule[\ruledef{s-trans}]
  {\Gamma\vdash\tau_1<:\tau_2
   \andalso
   \Gamma\vdash\tau_2<:\tau_3}
  {\Gamma\vdash\tau_1<:\tau_3}

\infax[\ruledef{s-bot}]
  {\Gamma\vdash\bot<:T}

\infax[\ruledef{s-top}]
  {\Gamma\vdash T<:\top}

\infrule[\ruledef{s-inter}]
  {\Gamma\vdash S<:T
   \andalso
   \Gamma\vdash S<:U}
  {\Gamma\vdash S<:\interty{T}{U}}

\infax[\ruledef{s-union-left}]
  {\Gamma\vdash T<:\unionty{T}{U}}

\infax[\ruledef{s-union-right}]
  {\Gamma\vdash U<:\unionty{T}{U}}

\infrule[\ruledef{s-union}]
  {\Gamma\vdash S<:U
   \andalso
   \Gamma\vdash T<:U}
  {\Gamma\vdash\unionty{S}{T}<:U}
\end{minipage}\hfill
\begin{minipage}[t]{0.51\linewidth}
\typicallabel{s-sel-hi}
\infrule[\ruledef{s-widen}]
  {\Gamma\vdash p\Rightarrow T}
  {\Gamma\vdash\sing{p}<:T}

\infrule[\ruledef{s-alias}]
  {\Gamma\vdash p\Rightarrow\sing{q}}
  {\Gamma\vdash\sing{q}<:\sing{p}}

\infrule[\ruledef{s-sel-hi}]
  {\Gamma\vdash p.A\Rightarrow S..T
   \andalso
   \Gamma\vdash S<:T}
  {\Gamma\vdash p.A<:T}

\infrule[\ruledef{s-sel-lo}]
  {\Gamma\vdash p.A\Rightarrow S..T
   \andalso
   \Gamma\vdash S<:T}
  {\Gamma\vdash S<:p.A}

\infax[\ruledef{s-inter-left}]
  {\Gamma\vdash\interty{T}{U}<:T}

\infax[\ruledef{s-inter-right}]
  {\Gamma\vdash\interty{T}{U}<:U}
\end{minipage}\par
\begin{minipage}[t]{\linewidth}
\typicallabel{s-bounds}
\infrule[\ruledef{s-fun}]
  {\Gamma\vdash S'<:S
   \andalso
   \Gamma,x:S'\vdash T<:T'}
  {\Gamma\vdash (x:S)\to T<:(x:S')\to T'}

\infrule[\ruledef{s-pair}]
  {\Gamma\vdash S<:S'
   \andalso
   \Gamma,x:S\vdash\tau<:\tau'}
  {\Gamma\vdash
    \pairty{x}{S}{a}{\tau}<:
    \pairty{x}{S'}{a}{\tau'}}

\infrule[\ruledef{s-merge}]
  {\Gamma\vdash S\mergeop T\Longrightarrow U}
  {\Gamma\vdash\interty{S}{T}<:U}

\infrule[\ruledef{s-bounds}]
  {\Gamma\vdash S'<:S
   \andalso
   \Gamma\vdash T<:T'
   \andalso
   \Gamma\vdash S<:T}
  {\Gamma\vdash S..T<:S'..T'}
\end{minipage}

\caption{Source syntax and static semantics (continued): signature
merge plans and subtyping.}
\Description{Both lower-join plans, all four aligned-merge plans, and all
eighteen signature-subtyping rules.}
\label{fig:subtyping}
\end{figure}

\begin{figure}[H]
\plateheader{Term typing}{\Gamma\vdash t:T}
\noindent
\begin{minipage}[t]{0.40\linewidth}
\typicallabel{t-path}
\infrule[\ruledef{t-path}]
  {\Gamma\vdash p\Rightarrow T}
  {\Gamma\vdash p:\sing{p}}

\infrule[\ruledef{t-abs}]
  {\Gamma,x:S\vdash t:T
   \andalso
   \Gamma\vdash S\ \wf}
  {\Gamma\vdash\lambda(x:S).t:(x:S)\to T}
\typicallabel{t-app}

\infrule[\ruledef{t-app}]
  {\Gamma\vdash p:(x:S)\to T
   \andalso
   \Gamma\vdash q:S}
  {\Gamma\vdash p\,q:\subst{T}{q}{x}}
\end{minipage}\hfill
\begin{minipage}[t]{0.57\linewidth}
\typicallabel{t-type-pair}
\infax[\ruledef{t-pair}]
  {\Gamma\vdash
    \pairval{y}{a}{z}:
    \pairty{x}{\sing{y}}{a}{\sing{z}}}

\infrule[\ruledef{t-type-pair}]
  {\Gamma\vdash T\ \wf}
  {\Gamma\vdash
    \pairval{y}{A}{T}:
    \pairty{x}{\sing{y}}{A}{T..T}}
\typicallabel{t-let}
\infrule[\ruledef{t-let}]
  {\Gamma\vdash t:S
   \andalso
   \Gamma\vdash T\ \wf
   \andalso
   \Gamma,x:S\vdash u:T}
  {\Gamma\vdash
    \mathbf{let}\ x=t\ \mathbf{in}\ u:T}

\infrule[\ruledef{t-sub}]
  {\Gamma\vdash t:S
   \andalso
   \Gamma\vdash S<:T
   \andalso
   \Gamma\vdash T\ \wf}
  {\Gamma\vdash t:T}
\end{minipage}

\par\addvspace{1.25em}
\noindent
\begin{minipage}[t]{0.57\linewidth}
\plateheader{Runtime syntax}{}
\[
\begin{array}{rcl}
r &::=& \loc{x} \mid \typeref{T}
\\
\sigma &::=& \varnothing
  \mid \sigma\mathbin{\cdot}(x\mapsto v)
\\
K &::=& [\,]
  \mid (\mathbf{let}\ x=[\,]\ \mathbf{in}\ u)::K
\\
c &::=& \langle\sigma,K,t\rangle
\end{array}
\]
\[
\referentof{z}=\loc{z}
\qquad
\referentof{T}=\typeref{T}.
\]
\end{minipage}\hfill
\begin{minipage}[t]{0.40\linewidth}
\plateheader{Final state}{\final(c)}
\typicallabel{f-loc}
\infax[\ruledef{f-loc}]
  {\final(\langle\sigma,[\,],x\rangle)}

\infax[\ruledef{f-val}]
  {\final(\langle\sigma,[\,],v\rangle)}
\end{minipage}

\caption{Term typing and run-time configurations: term typing, run-time
syntax, and final configurations.}
\Description{All seven term-typing rules, the complete run-time grammar,
and both final-state rules.}
\label{fig:term-typing}
\label{fig:runtime-syntax}
\end{figure}

\clearpage
\begin{figure}[H]
\ContinuedFloat

\begin{minipage}[t]{\linewidth}
\plateheader{Path resolution}{\sigma\vdash p\resolve r}
\typicallabel{res-miss}
\noindent
\begin{minipage}[t]{0.47\linewidth}
\infax[\ruledef{res-var}]
  {\sigma\vdash x\resolve\loc{x}}

\infrule[\ruledef{res-fst}]
  {\sigma\vdash p\resolve\loc{x}
   \andalso
   \sigma(x)=\pairval{y}{a}{\delta}}
  {\sigma\vdash p.1\resolve\loc{y}}
\end{minipage}\hfill
\begin{minipage}[t]{0.51\linewidth}
\infrule[\ruledef{res-sel}]
  {\sigma\vdash p\resolve\loc{x}
   \andalso
   \sigma(x)=\pairval{y}{a}{\delta}}
  {\sigma\vdash p.a\resolve\referentof{\delta}}

\infrule[\ruledef{res-miss}]
  {\sigma\vdash p\resolve\loc{x}
   \andalso
   \sigma(x)=\pairval{y}{b}{\delta}
   \\
   a\ne b
   \andalso
   \sigma\vdash y.a\resolve r}
  {\sigma\vdash p.a\resolve r}
\end{minipage}
\end{minipage}

\plateheader{Machine transition}
  {\langle\sigma,K,t\rangle\sstep
   \langle\sigma',K',t'\rangle}
\noindent
\begin{minipage}[t]{0.43\linewidth}
\typicallabel{e-app}
\infrule[\ruledef{e-app}]
   {\sigma\vdash p\resolve\loc{x}
   \andalso
   \sigma\vdash q\resolve\loc{y}
   \\
   \sigma(x)=\lambda(z:S).t}
  {\langle\sigma,K,p\,q\rangle
   \sstep\\
   \langle\sigma,K,\subst{t}{y}{z}\rangle}

\infrule[\ruledef{e-path}]
   {\sigma\vdash p\resolve\loc{x}
   \andalso
   p\ \mathsf{not\ a\ variable}}
  {\langle\sigma,K,p\rangle
   \sstep\\
   \langle\sigma,K,x\rangle}
\end{minipage}\hfill
\begin{minipage}[t]{0.54\linewidth}
\typicallabel{e-return}
\infax[\ruledef{e-let}]
  {\langle\sigma,K,\mathbf{let}\ x=t\ \mathbf{in}\ u\rangle
   \sstep\\
   \langle\sigma,
     (\mathbf{let}\ x=[\,]\ \mathbf{in}\ u)::K,t\rangle}

\infax[\ruledef{e-return}]
  {\langle\sigma,
      (\mathbf{let}\ x=[\,]\ \mathbf{in}\ u)::K,y\rangle
   \sstep\\
   \langle\sigma,K,\subst{u}{y}{x}\rangle}
\end{minipage}\par
\begin{minipage}[t]{\linewidth}
\typicallabel{e-alloc}
\infrule[\ruledef{e-alloc}]
  {z\ \mathsf{fresh}}
  {\langle\sigma,
      (\mathbf{let}\ x=[\,]\ \mathbf{in}\ u)::K,v\rangle
   \sstep\\
   \langle
     \sigma\mathbin{\cdot}(z\mapsto v),K,
     \subst{u}{z}{x}\rangle}
\end{minipage}

\caption{Term typing and run-time configurations (continued):
path resolution and CK-machine transitions.}
\Description{All four path-resolution rules and all five machine
transitions.}
\label{fig:path-resolution}
\label{fig:machine}
\end{figure}

\endgroup
\clearpage

We write \(\mathsf{initial}(t)=\langle\varnothing,[\,],t\rangle\) for the
initial state of a closed term \(t\).

\begin{lemma}[Resolution determinism]
\label{lem:resolution-determinism}
If \(\sigma\vdash p\resolve r\) and
\(\sigma\vdash p\resolve r'\), then \(r=r'\).
\end{lemma}

The proof is induction on the first resolution derivation. In the
nonmatching-selection case, both derivations follow the same store binding
and the induction hypothesis applies to the tail resolution.

\section{Challenges for preservation}
\label{sec:obligations}

\subsection{Singleton types and abstract type selections}

The distinction between singleton types and abstract type selections is required
by the alias rule. A premise \(\Gamma\vdash p\Rightarrow\sing{q}\) says that \(p\)
has the singleton type of \(q\), and therefore justifies reversing the alias:
\(\sing{q}<:\sing{p}\). A premise \(\Gamma\vdash p\Rightarrow q.A\) only places the
type of \(p\) between the bounds declared by \(q.A\).

An earlier syntax used one type constructor for both roles. Singleton
symmetry could then consume a premise assigning a path a selected type. With
\[
  f=\lambda(x:\top).x,
  \qquad q=\pairval{f}{A}{\top},
  \qquad h=\lambda(z:q.A).q,
\]
the conflated rules admitted the chain
\[
  \sing{q}<:\top<:q.A<:\sing{z}.
\]
They consequently assigned \(\sing{f}\) to \(h\,f\), although the
application returns \(q\). The complete closed program is
\[
\begin{array}{l}
\mathbf{let}\ f=\lambda(x:\top).x\ \mathbf{in}\\
\mathbf{let}\ q=\pairval{f}{A}{\top}\ \mathbf{in}\\
\mathbf{let}\ h=\lambda(z:q.A).q\ \mathbf{in}\\
\mathbf{let}\ g=h\,f\ \mathbf{in}\ g\,f.
\end{array}
\]
Evaluation binds \(g\) to the pair \(q\), while the erroneous singleton
type makes \(g\) usable as the function \(f\). The final application is
stuck because its operator cell contains a pair. Separate constructors
prevent a selected-type premise from serving as evidence of path identity. All intervals
in the example may have equal bounds, so the bound-consistency premises do
not address this problem.

\subsection{The dependent premise of pair subtyping}

The second premise of \ruleref{s-pair} is a derivation under a hypothetical
variable. An assignment of concrete locations to the variables in \(\Gamma\)
contains no location for that variable. Interpreting the premise immediately
would require choosing one before the source pair is observed. Ordinary
substitution lemmas also leave the critical semantic fact unstated: the
chosen location must have type \(S\), and the member referent must have the
source member signature opened at the same location.

A pair-realization derivation contains exactly these facts. If
\(\sigma(z)=\pairval{y}{a}{\delta}\) realizes
\(\pairty{x}{S}{a}{\tau}\), then \(y\) realizes \(S\) and the referent stored
by \(\delta\) realizes \(\tau[y/x]\). The coercion interpretation retains the
member-subtyping derivation until these realization derivations become
available, permitting \ruleref{s-pair} to vary both the first component and the member
signature.

\section{Store-indexed semantic interpretation}
\label{sec:coercions}

The proof has three stages. Typed path resolution interprets a path-synthesis
derivation as a run-time resolution together with a realization of its
referent. The subtyping interpretation compiles subtyping derivations to
store-indexed coercions, whose application maps source realizations to target
realizations. Premises beneath function and pair binders remain suspended:
function application supplies the argument location, while pair-realization
inversion supplies the stored first component. Finally, interpretation of
source typing produces syntax-directed typing for the running term;
continuation typing completes the invariant used for progress and
preservation (\cref{sec:invariant}).

\subsection{Store-induced path equality and signature conversion}

Path resolution may end at a location or at a stored type definition. We
write \(p\equiv_\sigma q\) for the least equivalence containing every pair
of paths that resolve to the same referent and closed under first projection
and member selection. Store-induced signature conversion
\(\tau\simeq_\sigma\tau'\) is the sort-preserving relation generated by
reflexivity, replacement of store-equal paths, and congruence for every
signature constructor:
\[
  p\equiv_\sigma q
  \quad\Longrightarrow\quad
  \subst{\tau_0}{p}{z}\simeq_\sigma\subst{\tau_0}{q}{z}.
\]
Replacement substitutes every free occurrence of \(z\), alpha-renaming an
intervening binder when necessary. Reversing the component path equalities
reverses a conversion. Sequences of conversions are handled by semantic
coercion composition. In particular, intersection and union conversion are
componentwise and neither reorder nor distribute their components.

Beneath a binder \(z\), the scoped path relation
\(\equiv_{\sigma;z}\) is the least path congruence containing
\(z\equiv_{\sigma;z}z\) and every equality inherited from
\(\equiv_\sigma\) between paths not containing \(z\). The corresponding
signature conversion \(\simeq_{\sigma;z}\) is generated by replacement with
this scoped path relation. It opens at any allocated location:
\[
  \tau\simeq_{\sigma;z}\tau'
  \quad y\in\mathsf{dom}(\sigma)
  \quad\Longrightarrow\quad
  \subst{\tau}{y}{z}\simeq_\sigma\subst{\tau'}{y}{z}.
\]
The complete unscoped rules appear in
\cref{fig:semantic-judgments}; scoped conversion uses the same rules with the
rigid path relation.

\subsection{Realization and coercions}

Fix a source context \(\Gamma\) and a store \(\sigma\). A closing
substitution \(\rho:\mathsf{dom}(\Gamma)\longrightarrow\mathsf{dom}(\sigma)\)
totally assigns allocated store locations to the source variables bound by
\(\Gamma\); distinct variables may share a location. For any path, type,
signature, member definition, or term \(X\), \(\rho(X)\) denotes
capture-avoiding application of \(\rho\), which replaces each free source
variable by its assigned location and leaves bound variables unchanged.
Extension \(\rho[z\mapsto y]\) maps the fresh variable \(z\) to
\(y\in\mathsf{dom}(\sigma)\) and satisfies
\[
  \subst{\rho(X)}{y}{z}=(\rho[z\mapsto y])(X)
\]
for any \(X\) formed beneath the binder \(z\).

The primary semantic judgments are
\[
\begin{array}{lll}
\rho\realizes_\sigma\Gamma
  &\quad&\text{the closing substitution \(\rho\) satisfies \(\Gamma\),}\\
\sigma\realizes x:T
  &&\text{location \(x\) realizes proper type \(T\),}\\
\sigma\realizes r:\tau
  &&\text{referent \(r\) realizes signature \(\tau\),}\\
\sigma\vdash\tau\coerce\tau'
  &&\text{a semantic coercion from \(\tau\) to \(\tau'\).}
\end{array}
\]
Realization is an inductive, store-indexed typing relation. A coercion is a
finite proof object whose application transforms a realization derivation at
a fixed referent. Capital letters \(Q\), \(R\), and \(C\) range over
resolution derivations, realization derivations, and coercions, respectively.
We write \(\mathcal R:\rho\realizes_\sigma\Gamma\) for a derivation of context
satisfaction.

Dependent premises use three auxiliary forms:
\[
\begin{array}{lll}
\bodycl{\sigma}{z}{S}{t}{T}
  &\quad&\text{a suspended body-typing derivation,}\\
\sigma;z:S\vdash T\coerce T'
  &&\text{an open codomain coercion,}\\
\membercl{\sigma}{z}{S}{\tau}{\tau'}
  &&\text{a suspended member-subtyping derivation.}
\end{array}
\]
The body and member forms suspend source derivations together with a closing
substitution satisfying their outer context. An open codomain coercion can be
composed or narrowed before the preservation proof instantiates it at an
argument location; \ruleref{cod-source} suspends the source subtyping derivation. A
suspended member premise is interpreted when inversion of a pair realization
supplies the stored first component. Alpha-equivalent choices of binder name
are identified.

Here \emph{closed} means that every free source variable has been replaced
by a location in \(\sigma\); a closed coercion may therefore mention store
locations. The judgments are mutually defined: realizations may contain
coercions, and coercions may contain realization derivations or suspended
source derivations. Their complete rules appear in
\cref{fig:semantic-judgments}.

Store-induced path equality preserves resolution: if
\(p\equiv_\sigma q\), then
\(\sigma\vdash p\resolve r\) exactly when
\(\sigma\vdash q\resolve r\). Store-induced signature conversion changes only
the paths occurring in a signature and is used solely in the metatheory.

\clearpage
\begingroup
\setlength{\abovedisplayskip}{5pt plus 1pt minus 1pt}
\setlength{\belowdisplayskip}{5pt plus 1pt minus 1pt}
\setlength{\intextsep}{4pt}
\captionsetup{skip=6pt}
\rulesetup
\footnotesize

\begin{figure}[H]
\setlength{\afterruleskip}{0.35em plus 0.08em minus 0.04em}
\noindent
\begin{minipage}[t]{0.48\linewidth}
\plateheader{Store-induced path equality}{p\equiv_\sigma q}
\typicallabel{peq-resolve}
\infax[\ruledef{peq-refl}]
  {p\equiv_\sigma p}

\infrule[\ruledef{peq-symm}]
  {p\equiv_\sigma q}
  {q\equiv_\sigma p}

\infrule[\ruledef{peq-trans}]
  {p\equiv_\sigma q
   \andalso
   q\equiv_\sigma p'}
  {p\equiv_\sigma p'}

\infrule[\ruledef{peq-resolve}]
  {\sigma\vdash p\resolve r
   \andalso
   \sigma\vdash q\resolve r}
  {p\equiv_\sigma q}

\infrule[\ruledef{peq-fst}]
  {p\equiv_\sigma q}
  {p.1\equiv_\sigma q.1}

\infrule[\ruledef{peq-sel}]
  {p\equiv_\sigma q}
  {p.a\equiv_\sigma q.a}
\end{minipage}\hfill
\begin{minipage}[t]{0.50\linewidth}
\plateheader{Store-induced signature conversion}{\tau\simeq_\sigma\tau'}
\typicallabel{seq-replace}
\infax[\ruledef{seq-refl}]
  {\tau\simeq_\sigma\tau}

\infrule[\ruledef{seq-replace}]
  {p\equiv_\sigma q}
  {\subst{\tau_0}{p}{z}\simeq_\sigma
   \subst{\tau_0}{q}{z}}

\infrule[\ruledef{seq-fun}]
  {S\simeq_\sigma S'
   \andalso
   T\simeq_{\sigma;z}T'}
  {(z:S)\to T\simeq_\sigma(z:S')\to T'}

\infrule[\ruledef{seq-pair}]
  {S\simeq_\sigma S'
   \andalso
   \tau\simeq_{\sigma;z}\tau'}
  {\pairty{z}{S}{a}{\tau}\simeq_\sigma
   \pairty{z}{S'}{a}{\tau'}}

\infrule[\ruledef{seq-inter}]
  {S_1\simeq_\sigma S_2
   \andalso
   T_1\simeq_\sigma T_2}
  {\interty{S_1}{T_1}\simeq_\sigma
   \interty{S_2}{T_2}}

\infrule[\ruledef{seq-union}]
  {S_1\simeq_\sigma S_2
   \andalso
   T_1\simeq_\sigma T_2}
  {\unionty{S_1}{T_1}\simeq_\sigma
   \unionty{S_2}{T_2}}

\infrule[\ruledef{seq-single}]
  {p\equiv_\sigma q}
  {\sing p\simeq_\sigma\sing q}

\infrule[\ruledef{seq-sel}]
  {p\equiv_\sigma q}
  {p.A\simeq_\sigma q.A}

\infrule[\ruledef{seq-intv}]
  {S_1\simeq_\sigma S_2
   \andalso
   T_1\simeq_\sigma T_2}
  {S_1..T_1\simeq_\sigma S_2..T_2}
\end{minipage}\par

\plateheader{Referent realization}{\sigma\realizes r:\tau}
\noindent
\begin{minipage}[t]{0.42\linewidth}
\typicallabel{ref-type}
\infrule[\ruledef{ref-loc}]
  {\sigma\realizes x:T}
  {\sigma\realizes\loc x:T}

\end{minipage}\hfill
\begin{minipage}[t]{0.55\linewidth}
\typicallabel{ref-type}
\infrule[\ruledef{ref-type}]
  {\sigma\vdash S\coerce T
   \andalso
   \sigma\vdash T\coerce U}
  {\sigma\realizes\typeref T:S..U}
\end{minipage}

\plateheader{Location realization}{\sigma\realizes x:T}
\noindent
\begin{minipage}[t]{0.50\linewidth}
\typicallabel{loc-single}
\infax[\ruledef{loc-top}]
  {\sigma\realizes x:\top}

\infrule[\ruledef{loc-inter}]
  {\sigma\realizes x:S
   \andalso
   \sigma\realizes x:T}
  {\sigma\realizes x:\interty{S}{T}}

\infrule[\ruledef{loc-union-left}]
  {\sigma\realizes x:S}
  {\sigma\realizes x:\unionty{S}{T}}

\infrule[\ruledef{loc-union-right}]
  {\sigma\realizes x:T}
  {\sigma\realizes x:\unionty{S}{T}}

\infrule[\ruledef{loc-fun}]
  {\sigma(x)=\lambda(z:S_0).t
   \quad
   \bodycl{\sigma}{z}{S_0}{t}{T_0}
   \\
   \sigma\vdash S\coerce S_0
   \quad
   \sigma;z:S\vdash T_0\coerce T}
  {\sigma\realizes x:(z:S)\to T}
\end{minipage}\hfill
\begin{minipage}[t]{0.48\linewidth}
\typicallabel{loc-select}
\infrule[\ruledef{loc-pair}]
  {\sigma(x)=\pairval{y}{a}{\delta}
   \andalso
   \sigma\realizes y:S
   \\
   \sigma\realizes\referentof{\delta}:\subst{\tau}{y}{z}}
  {\sigma\realizes x:\pairty{z}{S}{a}{\tau}}

\infrule[\ruledef{loc-single}]
  {\sigma\vdash p\resolve\loc x}
  {\sigma\realizes x:\sing p}

\infrule[\ruledef{loc-select}]
  {\sigma\vdash p.A\resolve\typeref T
   \andalso
   \sigma\realizes x:T}
  {\sigma\realizes x:p.A}
\end{minipage}
\caption{Store-indexed judgments (part 1): path equality, signature
conversion, and realization.}
\Description{All six path-equality rules, all nine signature-conversion rules,
both referent rules, and all eight location-realization rules.}
\label{fig:runtime-eq}
\label{fig:semantic-judgments}
\end{figure}

\clearpage
\begin{figure}[H]
\ContinuedFloat
\setlength{\afterruleskip}{0.09em plus 0.03em minus 0.03em}
\noindent
\begin{minipage}[t]{0.48\linewidth}
\plateheader{Context satisfaction}{\rho\realizes_\sigma\Gamma}
\typicallabel{env}
\infrule[\ruledef{env}]
  {\forall x\in\mathsf{dom}(\Gamma).\;
   \sigma\realizes\rho(x):\rho(\Gamma(x))}
  {\rho\realizes_\sigma\Gamma}

\end{minipage}\hfill
\begin{minipage}[t]{0.50\linewidth}
\plateheader{Suspended body typing}
  {\bodycl{\sigma}{z}{S}{t}{T}}
\typicallabel{body-source}
\infrule[\ruledef{body-source}]
  {\mathcal R:\rho\realizes_{\sigma}\Gamma
   \andalso
   \Gamma,z:S\vdash t:T}
  {\bodycl{\sigma}{z}{\rho(S)}{\rho(t)}{\rho(T)}}

\end{minipage}

\plateheader{Suspended member premise}
  {\membercl{\sigma}{z}{S}{\tau}{\tau'}}
\typicallabel{mem-source}
\infrule[\ruledef{mem-source}]
  {\mathcal R:\rho\realizes_\sigma\Gamma
   \andalso
   \Gamma,z:S\vdash\tau<:\tau'}
  {\membercl{\sigma}{z}{\rho(S)}{\rho(\tau)}{\rho(\tau')}}

\plateheader{Semantic coercion}
  {\sigma\vdash \tau_1\coerce \tau_2}
\noindent
\begin{minipage}[t]{0.43\linewidth}
\typicallabel{c-trans}
\infax[\ruledef{c-refl}]
  {\sigma\vdash \tau\coerce \tau}

\infrule[\ruledef{c-trans}]
  {\sigma\vdash \tau_1\coerce \tau_2
   \andalso
   \sigma\vdash \tau_2\coerce \tau_3}
  {\sigma\vdash \tau_1\coerce \tau_3}

\infrule[\ruledef{c-eq}]
  {\tau_1\simeq_\sigma\tau_2}
  {\sigma\vdash \tau_1\coerce \tau_2}

\infax[\ruledef{c-bot}]
  {\sigma\vdash\bot\coerce T}

\infax[\ruledef{c-top}]
  {\sigma\vdash T\coerce\top}

\infrule[\ruledef{c-bounds}]
  {\sigma\vdash S'\coerce S
   \andalso
   \sigma\vdash T\coerce T'}
  {\sigma\vdash S..T\coerce S'..T'}

\end{minipage}\hfill
\begin{minipage}[t]{0.54\linewidth}
\typicallabel{c-sel-hi}
\infrule[\ruledef{c-widen}]
  {\sigma\vdash p\resolve\loc x
   \andalso
   \sigma\realizes x:T}
  {\sigma\vdash\sing p\coerce T}

\infrule[\ruledef{c-alias}]
  {\sigma\vdash p\resolve\loc x
   \andalso
   \sigma\vdash q\resolve\loc x}
  {\sigma\vdash\sing q\coerce\sing p}

\infrule[\ruledef{c-sel-lo}]
  {\sigma\vdash p.A\resolve\typeref T
   \andalso
   \sigma\vdash S\coerce T}
  {\sigma\vdash S\coerce p.A}

\infrule[\ruledef{c-sel-hi}]
  {\sigma\vdash p.A\resolve\typeref T
   \andalso
   \sigma\vdash T\coerce U}
  {\sigma\vdash p.A\coerce U}

\infrule[\ruledef{c-fun}]
  {\sigma\vdash S'\coerce S
   \andalso
   \sigma;z:S'\vdash T\coerce T'}
  {\sigma\vdash
    (z:S)\to T\coerce(z:S')\to T'}

\infrule[\ruledef{c-pair}]
  {\sigma\vdash S\coerce S'
   \andalso
   \membercl{\sigma}{z}{S}{\tau}{\tau'}}
  {\sigma\vdash
    \pairty{z}{S}{a}{\tau}\coerce
    \pairty{z}{S'}{a}{\tau'}}
\end{minipage}

\typicallabel{c-merge}
\infrule[\ruledef{c-merge}]
  {\mathcal R:\rho\realizes_\sigma\Gamma
   \andalso
   \Gamma\vdash S\mergeop T\Longrightarrow U}
  {\sigma\vdash
   \interty{\rho(S)}{\rho(T)}\coerce\rho(U)}

\noindent
\begin{minipage}[t]{0.40\linewidth}
\typicallabel{c-inter}
\infrule[\ruledef{c-inter}]
  {\sigma\vdash R\coerce S
   \andalso
   \sigma\vdash R\coerce T}
  {\sigma\vdash R\coerce\interty{S}{T}}
\end{minipage}\hfill
\begin{minipage}[t]{0.28\linewidth}
\typicallabel{c-inter-left}
\infax[\ruledef{c-inter-left}]
  {\sigma\vdash\interty{S}{T}\coerce S}
\end{minipage}\hfill
\begin{minipage}[t]{0.28\linewidth}
\typicallabel{c-inter-right}
\infax[\ruledef{c-inter-right}]
  {\sigma\vdash\interty{S}{T}\coerce T}
\end{minipage}

\noindent
\begin{minipage}[t]{0.29\linewidth}
\typicallabel{c-union-left}
\infax[\ruledef{c-union-left}]
  {\sigma\vdash S\coerce\unionty{S}{T}}
\end{minipage}\hfill
\begin{minipage}[t]{0.29\linewidth}
\typicallabel{c-union-right}
\infax[\ruledef{c-union-right}]
  {\sigma\vdash T\coerce\unionty{S}{T}}
\end{minipage}\hfill
\begin{minipage}[t]{0.38\linewidth}
\typicallabel{c-union}
\infrule[\ruledef{c-union}]
  {\sigma\vdash S\coerce U
   \andalso
   \sigma\vdash T\coerce U}
  {\sigma\vdash\unionty{S}{T}\coerce U}
\end{minipage}

\plateheader{Open codomain coercion}
  {\sigma;z:S\vdash T\coerce T'}
\noindent
\begin{minipage}[t]{0.48\linewidth}
\typicallabel{cod-trans}
\infax[\ruledef{cod-refl}]
  {\sigma;z:S\vdash T\coerce T}

\infrule[\ruledef{cod-trans}]
  {\sigma;z:S\vdash T_1\coerce T_2
   \quad
   \sigma;z:S\vdash T_2\coerce T_3}
  {\sigma;z:S\vdash T_1\coerce T_3}

\infrule[\ruledef{cod-eq}]
  {T\simeq_{\sigma;z}T'}
  {\sigma;z:S\vdash T\coerce T'}
\end{minipage}\hfill
\begin{minipage}[t]{0.51\linewidth}
\typicallabel{cod-narrow}
\infrule[\ruledef{cod-narrow}]
  {\sigma\vdash S'\coerce S
   \andalso
   \sigma;z:S\vdash T\coerce T'}
  {\sigma;z:S'\vdash T\coerce T'}

\infrule[\ruledef{cod-source}]
  {\mathcal R:\rho\realizes_\sigma\Gamma
   \andalso
   \Gamma,z:S\vdash T<:T'}
  {\sigma;z:\rho(S)\vdash\rho(T)\coerce\rho(T')}
\end{minipage}
\caption{Store-indexed judgments (part 2): context satisfaction, suspended
premises, semantic coercions, and open codomain coercions.}
\Description{The context rule, the suspended-body rule, the suspended-member
rule, all nineteen semantic-coercion rules, and all five open-codomain rules.}
\end{figure}

\endgroup

We write \(\ruleuse{ref-loc}(R_x)\),
\(\ruleuse{ref-type}(C_L,C_U)\), and
\(\ruleuse{loc-single}(Q)\) for the derivations obtained by the corresponding
realization rules.

We abbreviate \(\ruleuse{c-trans}(C_1,C_2)\) by \(C_1;C_2\); the notation is
also used for \ruleref{cod-trans} on open codomain coercions.

Rule \ruleref{env} makes context satisfaction pointwise. Thus
\(\rho[z\mapsto y]\realizes_\sigma\Gamma,z:S\) whenever
\(\rho\realizes_\sigma\Gamma\) and
\(\sigma\realizes y:\rho(S)\). Every location names a store cell, so
\ruleref{loc-top} needs no separate lookup premise, and \(\bot\) has no
realization rule. Rule \ruleref{loc-inter} records two realizations of the
same location, one for each component; it imposes no condition on the stored
value beyond those already imposed by its premises. Rules
\ruleref{loc-union-left} and \ruleref{loc-union-right} instead record one
realization together with the tag identifying its union arm. This tag is
semantic evidence, not a run-time injection. Rule \ruleref{loc-fun} records the stored abstraction and its
body, a coercion from the ascribed argument type to the stored domain, and an
open codomain coercion from the stored body type to the ascribed result type.
Rule \ruleref{loc-pair} records realization of the first component and of the
member referent opened at that component.

A closed coercion contains the derivations needed to preserve realization at
the same referent. Identity and transitivity give identity and composition.
Rule \ruleref{c-eq} embeds store-induced signature conversion, whose soundness
follows from realization invariance below.
Singleton and selection coercions retain concrete resolution derivations.
Interval coercions retain the comparisons needed to reconstruct the target
bounds.
Rule \ruleref{c-inter} retains both coercions from the common source, while
\ruleref{c-inter-left} and \ruleref{c-inter-right} contain no additional
evidence: their action simply projects the corresponding realization.
The union injections add the corresponding semantic tag, while
\ruleref{c-union} retains both arm coercions and dispatches on that tag.
Rule \ruleref{c-merge} closes a source merge plan with the environment in
which it was compiled. Its structural action combines two realizations of
the same referent according to that plan. At a pair node it inverts both pair
realizations and uses functionality of store lookup to identify their one
stored pair, recursively merges the first-component evidence, extends the
saved environment with that concrete merged component, and only then
recursively merges the member evidence. At an interval node it combines the
lower coercions according to the lower-join plan and combines both upper
coercions before continuing with the upper merge plan. This single semantic
form covers the four former one-cell merge coercions and arbitrary aligned
record spines.

The dependent cases suspend their binder premises, as specified by
\ruleref{cod-source} and \ruleref{mem-source}. The two forms differ because an
open codomain coercion must survive further function subtyping, and may be
composed, transported by store-induced conversion, or narrowed before the
preservation proof supplies an argument location, whereas a suspended member
premise is interpreted immediately after pair-realization inversion.

\begin{lemma}[Constructor inversion for signature conversion]
\label{lem:conversion-structure}
If \(\tau\simeq_\sigma\tau'\), then the two signatures have the same outer
constructor. After alpha-renaming binders to a common name, the non-nullary
cases satisfy
\[
\begin{array}{rcl}
(z:S_1)\to T_1\simeq_\sigma(z:S_2)\to T_2
&\Longrightarrow&
S_1\simeq_\sigma S_2\ \land\ T_1\simeq_{\sigma;z}T_2,
\\[2pt]
\interty{S_1}{T_1}\simeq_\sigma\interty{S_2}{T_2}
&\Longrightarrow&
S_1\simeq_\sigma S_2\ \land\ T_1\simeq_\sigma T_2,
\\[2pt]
\unionty{S_1}{T_1}\simeq_\sigma\unionty{S_2}{T_2}
&\Longrightarrow&
S_1\simeq_\sigma S_2\ \land\ T_1\simeq_\sigma T_2,
\\[2pt]
\pairty{z}{S_1}{a}{\tau_1}\simeq_\sigma
\pairty{z}{S_2}{a'}{\tau_2}
&\Longrightarrow&
a=a'\ \land\ S_1\simeq_\sigma S_2\ \land\
\tau_1\simeq_{\sigma;z}\tau_2,
\\[2pt]
\sing p\simeq_\sigma\sing q
&\Longrightarrow& p\equiv_\sigma q,
\\[2pt]
p.A\simeq_\sigma q.A'
&\Longrightarrow& A=A'\ \land\ p\equiv_\sigma q,
\\[2pt]
S_1..T_1\simeq_\sigma S_2..T_2
&\Longrightarrow&
S_1\simeq_\sigma S_2\ \land\ T_1\simeq_\sigma T_2.
\end{array}
\]
In particular, \(\top\) and \(\bot\) convert only to themselves, and a
proper type never converts to an interval.
\end{lemma}

The proof is induction on signature conversion. Replacement preserves the
outer constructor by induction on the signature \(\tau_0\) in the replacement
rule; the congruence cases are immediate.

\begin{lemma}[Realization invariance]
\label{lem:conversion-transport}
If \(\tau\simeq_\sigma\tau'\) and
\(\sigma\realizes r:\tau\), then
\(\sigma\realizes r:\tau'\).
\end{lemma}

The proof is a simultaneous induction on referent and location realization,
using constructor inversion. Singleton and selection cases use preservation
of resolution under path equality. Function domains use conversion in the
contravariant direction; the open codomain coercion contained in the
function-realization derivation is narrowed to the target domain and composed
with the scoped codomain conversion. The intersection case transports both
component realizations independently and rebuilds \ruleref{loc-inter}; a
union case transports only its tagged arm and preserves the tag. In the pair
case, scoped conversion of member
signatures is opened at the stored first component before applying the
induction hypothesis to the member referent. Interval bounds are handled by
coercions induced by their component conversions.

\subsection{Interpreting source derivations}
\label{sec:coercion-construction}

Let \(\mathcal R:\rho\realizes_\sigma\Gamma\). Metavariables
\(\pi_p\) and \(\pi_s\) range over derivation trees of path synthesis and
subtyping, respectively. They are metatheoretic
objects and do not extend the source syntax or its judgments. Typed path
resolution and subtyping interpretation have signatures
\[
\begin{array}{rcll}
\resolvepi{\mathcal R}{\pi_p}=\langle r,Q,R\rangle
&\quad&(\pi_p\text{ derives }\Gamma\vdash p\Rightarrow\tau),
&
Q:\sigma\vdash\rho(p)\resolve r,
\\[2pt]
&&&R:\sigma\realizes r:\rho(\tau),
\\[2pt]
\compilepi{\mathcal R}{\pi_s}=C
&&(\pi_s\text{ derives }\Gamma\vdash\tau<:\tau'),
&C:\sigma\vdash\rho(\tau)\coerce\rho(\tau').
\end{array}
\]
The result depends on the chosen derivation: two derivations of the same
subtyping judgment may produce different coercions. Bracket notation denotes
this derivation-directed interpretation. The equations below define the
subtyping interpretation by its last source rule.
The rule label in the left margin identifies that rule, and \(\pi_i\) denotes its
immediate premise derivations in their displayed order in
\cref{fig:path-typing,fig:subtyping}.

Premises beneath a binder cannot be interpreted until a store location is
available. Let \(\pi_b\) derive \(\Gamma,z:S\vdash t:T\), let
\(\pi_m\) derive \(\Gamma,z:S\vdash\tau<:\tau'\), and let \(\pi_c\)
derive \(\Gamma,z:S\vdash T<:T'\). Suspending a premise pairs
\(\mathcal R\) with the corresponding open source derivation:
\[
\begin{array}{rcl}
\closure{\mathcal R}{\pi_b}{body}
  &:& \bodycl{\sigma}{z}{\rho(S)}{\rho(t)}{\rho(T)},
\\
\closure{\mathcal R}{\pi_m}{mem}
  &:& \membercl{\sigma}{z}{\rho(S)}{\rho(\tau)}{\rho(\tau')},
\\
\closure{\mathcal R}{\pi_c}{cod}
  &:& \sigma;z:\rho(S)\vdash\rho(T)\coerce\rho(T').
\end{array}
\]
The first line uses \ruleref{body-source}, the second \ruleref{mem-source}, and
the third \ruleref{cod-source}. Their types record whether the suspended
derivation is body typing, member subtyping, or codomain subtyping.

\begin{lemma}[Typed path resolution]
\label{lem:typed-resolution}
If \(\mathcal R:\rho\realizes_\sigma\Gamma\) and
\(\pi_p\) derives \(\Gamma\vdash p\Rightarrow\tau\), then
\[
  \resolvepi{\mathcal R}{\pi_p}=\langle r,Q,R\rangle
\]
for some referent \(r\), resolution derivation
\(Q:\sigma\vdash\rho(p)\resolve r\), and realization derivation
\(R:\sigma\realizes r:\rho(\tau)\).
\end{lemma}

The proof is induction on path synthesis. The two selection cases carry
the substance. For a matching label, pair realization supplies the stored
first component \(y\) and a realization of the member signature opened at
\(y\). The paths \(y\) and \(\rho(p).1\) resolve to the same location, so
store-induced path equality and \cref{lem:conversion-transport} transfer the
member realization to the signature synthesized by \ruleref{p-sel}. For a
nonmatching label, the induction hypothesis resolves the tail from
\(\rho(p).1\); equality with \(y\) transports that tail resolution before
\ruleref{res-miss} continues the lookup.

For \(R_y:\sigma\realizes y:\rho(S)\), the notation
\(\mathcal R[z\mapsto(y,R_y)]\) denotes the context-satisfaction derivation
\(\rho[z\mapsto y]\realizes_\sigma\Gamma,z:S\).

\begin{figure}[H]
\begingroup
\raggedright
\rulesetup
\footnotesize
\setlength{\afterruleskip}{0.45em plus 0.08em minus 0.05em}
\plateheader{Interpretation of subtyping derivations}
  {\compilepi{\mathcal R}{\pi}=C}
The bound-consistency premises of the selection rules produce no coercion.
Likewise, the third premise of \ruleref{s-bounds} establishes consistency of
the source interval but contributes no evidence to the resulting coercion.

\noindent
\begin{minipage}[t]{0.485\linewidth}
\typicallabel{s-trans}
\infax[\ruleuse{s-refl}]
  {\compilepi{\mathcal R}{\pi}=\ruleuse{c-refl}}

\infrule[\ruleuse{s-trans}]
  {\compilepi{\mathcal R}{\pi_1}=C_1
   \andalso
   \compilepi{\mathcal R}{\pi_2}=C_2}
  {\compilepi{\mathcal R}{\pi}=C_1;C_2}

\infax[\ruleuse{s-bot}]
  {\compilepi{\mathcal R}{\pi}=\ruleuse{c-bot}}

\infax[\ruleuse{s-top}]
  {\compilepi{\mathcal R}{\pi}=\ruleuse{c-top}}

\typicallabel{s-merge}
\infrule[\ruleuse{s-merge}]
  {\mathcal M:
    \Gamma\vdash S\mergeop T\Longrightarrow U}
  {\compilepi{\mathcal R}{\pi}=
    \ruleuse{c-merge}(\mathcal R,\mathcal M)}
\end{minipage}\hfill
\begin{minipage}[t]{0.515\linewidth}
\typicallabel{s-sel-hi}
\infrule[\ruleuse{s-widen}]
  {\begin{gathered}
   \resolvepi{\mathcal R}{\pi_1}=\\
   \langle\loc x,Q,\ruleuse{ref-loc}(R_T)\rangle
   \end{gathered}}
  {\compilepi{\mathcal R}{\pi}=\ruleuse{c-widen}(Q,R_T)}

\infrule[\ruleuse{s-alias}]
  {\begin{gathered}
   \resolvepi{\mathcal R}{\pi_1}=\\
   \langle\loc x,Q_p,
      \ruleuse{ref-loc}(\ruleuse{loc-single}(Q_q))\rangle
   \end{gathered}}
  {\compilepi{\mathcal R}{\pi}=\ruleuse{c-alias}(Q_p,Q_q)}

\infrule[\ruleuse{s-sel-hi}]
  {\begin{gathered}
   \resolvepi{\mathcal R}{\pi_1}=\\
   \langle\typeref W,Q,\ruleuse{ref-type}(C_L,C_U)\rangle
   \end{gathered}}
  {\compilepi{\mathcal R}{\pi}=\ruleuse{c-sel-hi}(Q,C_U)}

\infrule[\ruleuse{s-sel-lo}]
  {\begin{gathered}
   \resolvepi{\mathcal R}{\pi_1}=\\
   \langle\typeref W,Q,\ruleuse{ref-type}(C_L,C_U)\rangle
   \end{gathered}}
  {\compilepi{\mathcal R}{\pi}=\ruleuse{c-sel-lo}(Q,C_L)}
\end{minipage}

\plateheader{Intersection cases}{}
\noindent
\begin{minipage}[t]{0.40\linewidth}
\typicallabel{s-inter}
\infrule[\ruleuse{s-inter}]
  {\compilepi{\mathcal R}{\pi_1}=C_1
   \andalso
   \compilepi{\mathcal R}{\pi_2}=C_2}
  {\compilepi{\mathcal R}{\pi}=\ruleuse{c-inter}(C_1,C_2)}
\end{minipage}\hfill
\begin{minipage}[t]{0.28\linewidth}
\typicallabel{s-inter-left}
\infax[\ruleuse{s-inter-left}]
  {\begin{gathered}
   \compilepi{\mathcal R}{\pi}=\\[-1pt]\ruleuse{c-inter-left}
   \end{gathered}}
\end{minipage}\hfill
\begin{minipage}[t]{0.28\linewidth}
\typicallabel{s-inter-right}
\infax[\ruleuse{s-inter-right}]
  {\begin{gathered}
   \compilepi{\mathcal R}{\pi}=\\[-1pt]\ruleuse{c-inter-right}
   \end{gathered}}
\end{minipage}

\plateheader{Union cases}{}
\noindent
\begin{minipage}[t]{0.29\linewidth}
\typicallabel{s-union-left}
\infax[\ruleuse{s-union-left}]
  {\begin{gathered}
   \compilepi{\mathcal R}{\pi}=\\[-1pt]\ruleuse{c-union-left}
   \end{gathered}}
\end{minipage}\hfill
\begin{minipage}[t]{0.29\linewidth}
\typicallabel{s-union-right}
\infax[\ruleuse{s-union-right}]
  {\begin{gathered}
   \compilepi{\mathcal R}{\pi}=\\[-1pt]\ruleuse{c-union-right}
   \end{gathered}}
\end{minipage}\hfill
\begin{minipage}[t]{0.38\linewidth}
\typicallabel{s-union}
\infrule[\ruleuse{s-union}]
  {\compilepi{\mathcal R}{\pi_1}=C_1
   \andalso
   \compilepi{\mathcal R}{\pi_2}=C_2}
  {\compilepi{\mathcal R}{\pi}=\ruleuse{c-union}(C_1,C_2)}
\end{minipage}

\plateheader{Dependent cases}{}

\infrule[\ruleuse{s-fun}]
  {\compilepi{\mathcal R}{\pi_1}=C_D
   \andalso F=\closure{\mathcal R}{\pi_2}{cod}}
  {\compilepi{\mathcal R}{\pi}=\ruleuse{c-fun}(C_D,F)}

\infrule[\ruleuse{s-pair}]
  {\compilepi{\mathcal R}{\pi_1}=C_S
   \andalso M=\closure{\mathcal R}{\pi_2}{mem}}
  {\compilepi{\mathcal R}{\pi}=\ruleuse{c-pair}(C_S,M)}

\infrule[\ruleuse{s-bounds}]
  {\compilepi{\mathcal R}{\pi_1}=C_1
   \andalso \compilepi{\mathcal R}{\pi_2}=C_2}
  {\compilepi{\mathcal R}{\pi}=\ruleuse{c-bounds}(C_1,C_2)}
\endgroup
\caption{Interpretation of all eighteen subtyping rules as semantic coercions.}
\Description{The eighteen equations defining subtyping interpretation. The
function and pair-covariance coercions retain their dependent premise
derivations; the merge coercion retains its source plan and environment.}
\label{fig:subtyping-interpretation}
\end{figure}

The dependent cases retain \(\pi_2\) together with \(\mathcal R\) instead of
interpreting it recursively. The closing substitution in \(\mathcal R\)
covers \(\Gamma\), but not the variable bound by the premise. For functions,
the retained derivation becomes an open codomain coercion that is instantiated
when application supplies the argument location. For pairs, it becomes a
suspended member premise that is instantiated after pair-realization inversion
supplies the stored first component. Extending the closing substitution with
that location then interprets \(\pi_2\) as a closed coercion.

\subsection{Coercion application}
\label{sec:coercion-action}

For \(C:\sigma\vdash\tau\coerce\tau'\) and
\(R:\sigma\realizes r:\tau\), coercion application has the signature
\[
  \action{C}{R}:\sigma\realizes r:\tau'.
\]
Application follows the coercion derivation. Reflexivity returns \(R\), and
composition applies its two coercions in sequence:
\[
  \action{\ruleuse{c-refl}}{R}=R,
  \qquad
  \action{C_1;C_2}{R}=\action{C_2}{\action{C_1}{R}}.
\]
For \ruleref{c-eq}, \cref{lem:conversion-transport} supplies the target
realization. Widening and aliasing use the resolution derivations retained in
their coercions; resolution determinism identifies their locations with the
referent in \(R\). The \ruleref{c-bot} case is vacuous because \(\bot\) has
no realization.

Intersection introduction runs both retained coercions on the same source
realization. If
\[
  \action{C_1}{\ruleuse{ref-loc}(R_0)}=\ruleuse{ref-loc}(R_1),
  \qquad
  \action{C_2}{\ruleuse{ref-loc}(R_0)}=\ruleuse{ref-loc}(R_2),
\]
then
\[
  \action{\ruleuse{c-inter}(C_1,C_2)}
    {\ruleuse{ref-loc}(R_0)}
  =\ruleuse{ref-loc}(\ruleuse{loc-inter}(R_1,R_2)).
\]
The elimination coercions merely select stored semantic evidence:
\[
\begin{array}{rcl}
  \action{\ruleuse{c-inter-left}}
    {\ruleuse{ref-loc}(\ruleuse{loc-inter}(R_1,R_2))}
  &=&\ruleuse{ref-loc}(R_1),
  \\[2pt]
  \action{\ruleuse{c-inter-right}}
    {\ruleuse{ref-loc}(\ruleuse{loc-inter}(R_1,R_2))}
  &=&\ruleuse{ref-loc}(R_2).
\end{array}
\]
This packaging exists only in realization evidence; it is not a source or
runtime term constructor.

Union coercions manipulate the tags of realization evidence. The two
injections wrap an existing realization without changing the location:
\[
\begin{array}{rcl}
  \action{\ruleuse{c-union-left}}{\ruleuse{ref-loc}(R_S)}
  &=&\ruleuse{ref-loc}(\ruleuse{loc-union-left}(R_S)),
  \\[2pt]
  \action{\ruleuse{c-union-right}}{\ruleuse{ref-loc}(R_T)}
  &=&\ruleuse{ref-loc}(\ruleuse{loc-union-right}(R_T)).
\end{array}
\]
Elimination dispatches to the coercion for the arm named by that tag:
\[
\begin{array}{rcl}
  \action{\ruleuse{c-union}(C_S,C_T)}
    {\ruleuse{ref-loc}(\ruleuse{loc-union-left}(R_S))}
  &=&\action{C_S}{\ruleuse{ref-loc}(R_S)},
  \\[2pt]
  \action{\ruleuse{c-union}(C_S,C_T)}
    {\ruleuse{ref-loc}(\ruleuse{loc-union-right}(R_T))}
  &=&\action{C_T}{\ruleuse{ref-loc}(R_T)}.
\end{array}
\]
The tag exists only in semantic evidence; source and run-time terms remain
untagged.

Merge coercions also act only on evidence. Let
\(\mathsf{merge}_{\mathcal R}(\mathcal M;R_1,R_2)\) denote the structural
action of a plan
\(\mathcal M:\Gamma\vdash\tau_1\mergeop\tau_2\Longrightarrow\tau\)
on two realizations of the same referent under \(\mathcal R\). Then
\[
\begin{aligned}
 &\action{\ruleuse{c-merge}(\mathcal R,\mathcal M)}
   {\ruleuse{ref-loc}(\ruleuse{loc-inter}(R_1,R_2))}
 \\
 &\qquad=
   \mathsf{merge}_{\mathcal R}
     (\mathcal M;\ruleuse{ref-loc}(R_1),\ruleuse{ref-loc}(R_2)).
\end{aligned}
\]
For a plan ending in \ruleref{m-same}, this action returns the first
realization. For a proper-type plan ending in \ruleref{m-inter}, it packages
the two location realizations with \ruleref{loc-inter}. These two choices
explain why merge normalization is plan directed: the same pair of equal or
aligned types can deliberately be kept opaque rather than recursively
normalized.

The pair case contains the binder-sensitive step. Suppose the two inputs
invert to
\[
\begin{array}{rcl}
  R_1&=&\ruleuse{ref-loc}
    (\ruleuse{loc-pair}(B_x,R_{S_1},R_{\tau_1})),\\[2pt]
  R_2&=&\ruleuse{ref-loc}
    (\ruleuse{loc-pair}(B_x',R_{S_2},R_{\tau_2})).
\end{array}
\]
Both bindings concern the same receiver location. Functionality of immutable
store lookup identifies the complete stored pair: in particular, both views
have the same stored first location \(y\), label \(a\), member definition
\(\delta\), and member referent. Let the pair plan have first subplan
\(\mathcal M_S\) and member subplan \(\mathcal M_\tau\). First compute
\[
  \ruleuse{ref-loc}(R_S)
  =\mathsf{merge}_{\mathcal R}
     (\mathcal M_S;
       \ruleuse{ref-loc}(R_{S_1}),
       \ruleuse{ref-loc}(R_{S_2})).
\]
If \(\mathcal M_S\) produces \(S\), then \(R_S\) proves that the concrete
location \(y\) realizes \(\rho(S)\). Extend the saved environment with that
fact:
\[
  \mathcal R'=\mathcal R[x\mapsto(y,R_S)]
  :\rho[x\mapsto y]\realizes_\sigma\Gamma,x:S.
\]
The member subplan was formed under exactly this merged binder. Both stored
member realizations have already been opened at the same \(y\), so after the
standard closing-substitution equation they can be merged as
\[
  R_\tau=
  \mathsf{merge}_{\mathcal R'}
    (\mathcal M_\tau;R_{\tau_1},R_{\tau_2}).
\]
The result is one realization of the unchanged physical cell:
\[
  \ruleuse{ref-loc}
    (\ruleuse{loc-pair}(B_x,R_S,R_\tau)).
\]
Thus recursive normalization may change the record prefix and a dependent
outer member simultaneously, while the environment extension gives the
member plan the concrete interpretation of its merged binder. Recursing on
the first subplan walks older cells in an aligned record spine. Requiring the
same label and member sort at every pair node is essential: one immutable
cell contains one label and one definition, and the calculus has no row
permutation with which to realign differently ordered fields.

For an interval plan, both realizations concern the same stored type witness
\(W\):
\[
  \ruleuse{ref-type}(C_{L_1},C_{U_1}),
  \qquad
  \ruleuse{ref-type}(C_{L_2},C_{U_2}).
\]
If the lower-join subplan ends in \ruleref{j-same}, either lower coercion is
reused. If it ends in \ruleref{j-union}, union elimination gives
\[
  \ruleuse{c-union}(C_{L_1},C_{L_2}):
  \sigma\vdash\unionty{L_1}{L_2}\coerce W.
\]
For an upper subplan
\(\mathcal M_U:\Gamma\vdash U_1\mergeop U_2\Longrightarrow U\), first use
intersection introduction and then the merge coercion retained for that
subplan:
\[
  \ruleuse{c-inter}(C_{U_1},C_{U_2});
  \ruleuse{c-merge}(\mathcal R,\mathcal M_U)
  :\sigma\vdash W\coerce U.
\]
Those lower and upper coercions reconstruct realization of the merged
interval. This action is semantically valid whether or not the target
interval is syntactically consistent. A source term using the result must
independently prove the target well formed, as required by \ruleref{t-sub}.

Selection coercions use the lower or upper coercion contained in the realized
interval. For lower selection, if applying
\(C_L:\sigma\vdash S\coerce W\) gives
\[
  \action{C_L}{\ruleuse{ref-loc}(R_x)}=
  \ruleuse{ref-loc}(R_W),
\]
then the retained resolution derivation
\(Q:\sigma\vdash p.A\resolve\typeref W\) reconstructs the selection
realization:
\[
  \action{\ruleuse{c-sel-lo}(Q,C_L)}
    {\ruleuse{ref-loc}(R_x)}
  =\ruleuse{ref-loc}(\ruleuse{loc-select}(Q,R_W)).
\]
For upper selection, inversion exposes the witness realization at the stored
type. Resolution determinism identifies the two stored types, after which the
upper coercion is applied:
\[
  \action{\ruleuse{c-sel-hi}(Q,C_U)}
    {\ruleuse{ref-loc}(\ruleuse{loc-select}(Q',R_W))}
  =\action{C_U}{\ruleuse{ref-loc}(R_W)},
\]
where determinism has identified the result type of \(Q'\) with \(W\).

The function and interval cases compose the component coercions. Suppose a
function realization contains a domain coercion
\(C_I:\sigma\vdash S\coerce S_0\) and an open codomain coercion
\(F_I:\sigma;z:S\vdash T_0\coerce T\). Applying
\(\ruleuse{c-fun}(C_D,F)\), where
\(C_D:\sigma\vdash S'\coerce S\) and
\(F:\sigma;z:S'\vdash T\coerce T'\), replaces them by
\[
  C_D;C_I:\sigma\vdash S'\coerce S_0,
  \qquad
  \ruleuse{cod-narrow}(C_D,F_I);F:
    \sigma;z:S'\vdash T_0\coerce T'.
\]
For an interval realization of \(\typeref W:S..T\), let
\(C_L:\sigma\vdash S\coerce W\) and
\(C_U:\sigma\vdash W\coerce T\). Applying
\(\ruleuse{c-bounds}(C_1,C_2)\), with
\(C_1:\sigma\vdash S'\coerce S\),
and \(C_2:\sigma\vdash T\coerce T'\), constructs the target realization
\[
  C_1;C_L:\sigma\vdash S'\coerce W,
  \qquad
  C_U;C_2:\sigma\vdash W\coerce T'.
\]
The consistency premise of the source subtyping rule is needed to admit the
interval, but coercion application uses only these two component coercions.

The \ruleref{c-pair} covariance case is the only application clause that
interprets a suspended source derivation. It is stated in full below.

\begin{figure}[H]
\begingroup
\raggedright
\hfuzz=2pt
\rulesetup
\plateheader{Member-premise instantiation}
  {\forcemem{M}{R_y}}
For \(M:\membercl{\sigma}{z}{S}{\tau}{\tau'}\) and
\(R_y:\sigma\realizes y:S\), instantiation has result type
\[
\forcemem{M}{R_y}:
\sigma\vdash\subst{\tau}{y}{z}\coerce\subst{\tau'}{y}{z}.
\]
Its sole equation resumes the exact derivation retained by
\ruleref{mem-source}:
\[
\forcemem{\closure{\mathcal R}{\pi_s}{mem}}{R_y}
=\compilepi{\mathcal R[z\mapsto(y,R_y)]}{\pi_s}.
\]

\plateheader{Open-codomain instantiation}
  {\forcecod{F}{R_y}}
For \(F:\sigma;z:S\vdash T\coerce T'\) and
\(R_y:\sigma\realizes y:S\), instantiation has result type
\[
\forcecod{F}{R_y}:
\sigma\vdash\subst{T}{y}{z}\coerce\subst{T'}{y}{z}.
\]

Instantiation preserves identity, composition, and store-induced conversion.
Narrowing first converts the argument realization to the domain expected by
the remaining open codomain coercion:
\[
\forcecod{\ruleuse{cod-narrow}(C_D,F)}{R_y}
  =\forcecod{F}{R_y'}
\quad\text{where}\quad
\action{C_D}{\ruleuse{ref-loc}(R_y)}=\ruleuse{ref-loc}(R_y').
\]
The \ruleref{cod-source} case resumes its suspended derivation:
\[
\forcecod{\closure{\mathcal R}{\pi_s}{cod}}{R_y}
  =\compilepi{\mathcal R[z\mapsto(y,R_y)]}{\pi_s}.
\]
\endgroup
\caption{Instantiation of suspended member and codomain premises.}
\Description{The result types of member-premise and open-codomain
instantiation, together with their source-derivation equations.}
\end{figure}

Coercion application preserves the referent; it changes only the realization
derivation, while the store and machine state remain fixed.

\subsection{The dependent-pair case}
\label{sec:coercion-pair}

Fix
\[
  \mathcal R:\rho\realizes_\sigma\Gamma,
  \qquad
  \pi_1:\Gamma\vdash S<:S',
  \qquad
  \pi_2:\Gamma,z:S\vdash\tau<:\tau'.
\]
For a derivation ending in \ruleref{s-pair}, interpretation yields
\[
  C_S=\compilepi{\mathcal R}{\pi_1},
  \qquad
  M=\closure{\mathcal R}{\pi_2}{mem},
  \qquad
  C=\ruleuse{c-pair}(C_S,M).
\]
Suppose
\[
  R:\sigma\realizes\loc x:
    \pairty{z}{\rho(S)}{a}{\rho(\tau)}.
\]
Inversion of \(R\) by \ruleref{ref-loc} and \ruleref{loc-pair} supplies
\[
  B_x:\sigma(x)=\pairval{y}{a}{\delta},
  \qquad
  R_y:\sigma\realizes y:\rho(S),
  \qquad
  R_\delta:\sigma\realizes\referentof{\delta}:
    \subst{\rho(\tau)}{y}{z}.
\]
Coercion application computes
\[
  \action{C_S}{\ruleuse{ref-loc}(R_y)}=\ruleuse{ref-loc}(R_y'),
  \qquad
  C_\tau=\forcemem{M}{R_y},
  \qquad
  R_\delta'=\action{C_\tau}{R_\delta}.
\]
The first equation gives
\(R_y':\sigma\realizes y:\rho(S')\). Member instantiation extends
\(\mathcal R\) with the original realization
\(R_y:\sigma\realizes y:\rho(S)\), resumes \(\pi_2\), and gives
\[
  C_\tau:\sigma\vdash
  \subst{\rho(\tau)}{y}{z}\coerce
  \subst{\rho(\tau')}{y}{z}.
\]
The last equation gives
\[
  R_\delta':\sigma\realizes\referentof{\delta}:
    \subst{\rho(\tau')}{y}{z}.
\]
Rule \ruleref{loc-pair}, with the unchanged binding \(B_x\), now derives
\[
  \sigma\realizes x:
    \pairty{z}{\rho(S')}{a}{\rho(\tau')},
\]
and \ruleref{ref-loc} supplies \(\action{C}{R}\). The first-component
location remains \(y\); only its realization and the member realization are
reconstructed at the target signatures. The same argument covers term and
type members, since \(\referentof{\delta}\) is respectively a location or a
stored type definition.

\subsection{Well-foundedness of coercion application}
\label{sec:coercion-termination}

The transitivity, intersection-introduction, union-elimination, and selection
cases recurse on coercions appearing as premises of the current coercion derivation. The
\ruleref{c-pair} covariance case also invokes the coercion obtained by
instantiating its suspended member premise. That coercion can be larger than
the enclosing pair coercion, so coercion structure alone does not justify
recursion. The allocation order of the immutable store provides the required
decrease.

Give a location referent \(\loc{x}\) its position in the allocation order,
starting at one, and give every type referent rank zero. Write this number as
\(\mathsf{rank}_\sigma(r)\). If
\[
  \sigma(x)=\pairval{y}{a}{\delta},
\]
then both referents stored in the pair have smaller rank:
\[
  \mathsf{rank}_\sigma(\loc y)
    <\mathsf{rank}_\sigma(\loc x),
  \qquad
  \mathsf{rank}_\sigma(\referentof{\delta})
    <\mathsf{rank}_\sigma(\loc x).
\]
For a type member, \(\referentof{\delta}=\typeref{T}\) has rank zero. For a
term member it is an older location named by the stored pair.
Fresh allocation preserves the relative rank of all existing referents.

Let \(|C|\) measure the structural size of a closed coercion, ignoring
resolution derivations and suspended premises. Coercion application is
defined by well-founded induction on
\[
  \bigl(\mathsf{rank}_\sigma(r),|C|\bigr)
\]
in lexicographic order. The premises of \ruleref{c-trans},
\ruleref{c-inter}, \ruleref{c-union}, \ruleref{c-sel-lo}, and
\ruleref{c-sel-hi} retain the referent and recurse on a structurally smaller
coercion, decreasing the secondary component. Union elimination inspects the
realization tag and recurses on only the corresponding, smaller arm coercion.
The intersection projections and union injections do not recurse. In
\ruleref{c-pair}, the first recursive application concerns \(\loc y\), and the
second concerns \(\referentof{\delta}\).  Both strictly decrease the primary
component, independently of the size of the coercion obtained by member
instantiation. Rule \ruleref{c-merge} also has coercion size one and invokes a
separate structural action on its retained finite merge plan. That action
recurses only on strict first-component and member subplans. At an interval
node it constructs, but does not immediately apply, the coercion retained for
the strict upper subplan. Consequently merge-plan action terminates by the
structural size of the plan and adds no recursive case to the lexicographic
coercion-action argument above. These are all recursive coercion applications.

\subsection{Semantic subtyping}
\label{sec:coercion-totality}

\begin{lemma}[Subtyping interpretation]
\label{lem:coercion-construction}
If \(\mathcal R:\rho\realizes_\sigma\Gamma\) and
\(\pi_s\) derives \(\Gamma\vdash\tau_1<:\tau_2\), then
\[
  \compilepi{\mathcal R}{\pi_s}:
  \sigma\vdash\rho(\tau_1)\coerce\rho(\tau_2).
\]
\end{lemma}

The proof is induction on the source derivation using the equations in
\cref{fig:subtyping-interpretation}. Typed path resolution supplies the
referent and its realization at the synthesized signature in the widening,
alias, and selection
cases. The dependent premise of a function or pair is suspended rather than
interpreted under a closing substitution that lacks its bound variable.

\begin{lemma}[Coercion application]
\label{lem:coercion-action}
If \(C:\sigma\vdash\tau_1\coerce\tau_2\) and
\(R:\sigma\realizes r:\tau_1\), then
\[
  \action{C}{R}:\sigma\realizes r:\tau_2.
\]
\end{lemma}

The referent remains \(r\); application changes only its realization
derivation. The well-founded induction needed for the pair-covariance case was established in
\cref{sec:coercion-termination}.

Subtyping interpretation followed by coercion application gives the semantic
content of a source subtyping derivation.

\begin{theorem}[Semantic subtyping]
\label{thm:semantic-subtyping}
If \(\mathcal R:\rho\realizes_\sigma\Gamma\),
\(\pi_s\) derives \(\Gamma\vdash\tau<:\tau'\), and
\(R:\sigma\realizes r:\rho(\tau)\), then
\[
  \action{\compilepi{\mathcal R}{\pi_s}}{R}:
  \sigma\realizes r:\rho(\tau').
\]
\end{theorem}

\section{The machine invariant}
\label{sec:invariant}

\subsection{Syntax-directed machine typing}

Source typing permits subsumption at any depth, so direct inversion may first
expose an arbitrary \ruleref{t-sub} derivation. The machine invariant uses
syntax-directed typing and places subsumption at the boundary of each value
or non-value term derivation. Its value and term judgments,
\(\sigma\vdash_{\mathsf v}v:T\) and
\(\sigma\vdash_{\mathsf{tm}}t:T\), relate the result type of the
corresponding syntax-directed source rule to \(T\) by a final coercion;
a machine term that is already a value delegates to value typing. These are
metatheoretic judgments over the run-time syntax.

\begin{figure}[H]
\begingroup
\raggedright
\rulesetup
\judgheader{Runtime value typing}{\sigma\vdash_{\mathsf v} v:T}
\noindent
\begin{minipage}[t]{\linewidth}
\typicallabel{val-abs}
\infrule[\ruledef{val-abs}]
  {\bodycl{\sigma}{x}{S_0}{t}{T_0}
   \andalso
   \sigma\vdash (x:S_0)\to T_0\coerce T}
  {\sigma\vdash_{\mathsf v}\lambda(x:S_0).t:T}
\end{minipage}\par
\begin{minipage}[t]{0.49\linewidth}
\typicallabel{val-pair}
\infrule[\ruledef{val-pair}]
  {\sigma\vdash
    \pairty{x}{\sing y}{a}{\sing z}\coerce T}
  {\sigma\vdash_{\mathsf v}\pairval{y}{a}{z}:T}
\end{minipage}\hfill
\begin{minipage}[t]{0.505\linewidth}
\typicallabel{val-type-pair}
\infrule[\ruledef{val-type-pair}]
  {\sigma\vdash
    \pairty{x}{\sing y}{A}{S..S}\coerce T}
  {\sigma\vdash_{\mathsf v}\pairval{y}{A}{S}:T}
\end{minipage}

\judgheader{Machine-term typing}{\sigma\vdash_{\mathsf{tm}} t:T}
\noindent
\begin{minipage}[t]{0.43\linewidth}
\typicallabel{term-value}
\infrule[\ruledef{term-path}]
  {\sigma\vdash p\resolve\loc x
   \andalso
   \sigma\vdash\sing p\coerce T}
  {\sigma\vdash_{\mathsf{tm}} p:T}

\infrule[\ruledef{term-value}]
  {\sigma\vdash_{\mathsf v} v:T}
  {\sigma\vdash_{\mathsf{tm}} v:T}
\end{minipage}\hfill
\begin{minipage}[t]{0.55\linewidth}
\typicallabel{term-let}
\infrule[\ruledef{term-app}]
  {\sigma\vdash_{\mathsf{tm}} p:(x:S)\to U
   \andalso
   \sigma\vdash_{\mathsf{tm}} q:S
   \\
   \sigma\vdash\subst{U}{q}{x}\coerce T}
  {\sigma\vdash_{\mathsf{tm}} p\,q:T}

\infrule[\ruledef{term-let}]
  {\sigma\vdash_{\mathsf{tm}} t:S
   \andalso
   \bodycl{\sigma}{x}{S}{u}{U}
   \\
   \sigma\vdash U\coerce T}
  {\sigma\vdash_{\mathsf{tm}}\mathbf{let}\ x=t\ \mathbf{in}\ u:T}
\end{minipage}

\judgheader{Continuation typing}{\sigma\vdash K:S\Rightarrow T}
\noindent
\begin{minipage}[t]{0.43\linewidth}
\typicallabel{kont-hole}
\infax[\ruledef{kont-hole}]
  {\sigma\vdash[\,]:T\Rightarrow T}
\end{minipage}\hfill
\begin{minipage}[t]{0.55\linewidth}
\typicallabel{kont-let}
\infrule[\ruledef{kont-let}]
  {\sigma\vdash K:U\Rightarrow T
   \andalso
   \bodycl{\sigma}{x}{S}{u}{U'}
   \\
   \sigma\vdash U'\coerce U}
  {\sigma\vdash
    (\mathbf{let}\ x=[\,]\ \mathbf{in}\ u)::K:S\Rightarrow T}
\end{minipage}

\judgheader{Configuration typing}{\vdash\langle\sigma,K,t\rangle:T}
\typicallabel{cfg-state}
\infrule[\ruledef{cfg-state}]
  {\sigma\vdash_{\mathsf{tm}} t:S
   \andalso
   \sigma\vdash K:S\Rightarrow T}
  {\vdash\langle\sigma,K,t\rangle:T}
\endgroup
\caption{The complete store-indexed machine invariant: syntax-directed
typing for values, terms, continuations, and configurations.}
\Description{All three value-typing rules, four machine-term rules, two
continuation-typing rules, and the configuration-typing rule.}
\label{fig:machine-typing}
\label{fig:configuration-typing}
\end{figure}

The two pair-value rules begin from the types assigned by the corresponding
source pair rules. In \ruleref{term-let}, \(U\) is independent of the let
binder. Continuations are typed as transformations from the type of the
current term to the final result type; the empty continuation is the
identity transformation.

For \(c=\langle\sigma,K,t\rangle\), write \(\vdash c:T\) for configuration
typing in \cref{fig:configuration-typing}.

\subsection{Interpretation of source typing}

Machine typing is closed under further subtyping. Value rules and non-value
term rules carry a final coercion from their syntax-directed result type to
the type stated by the judgment; another coercion is incorporated by
composition. The value-delegation rule uses the corresponding value result.

\begin{lemma}[Machine subsumption]
\label{lem:machine-postcomposition}
If \(\sigma\vdash_{\mathsf{tm}} t:S\) and \(C:\sigma\vdash S\coerce T\), then
\(\sigma\vdash_{\mathsf{tm}} t:T\). Likewise, if
\(\sigma\vdash_{\mathsf v} v:S\) and \(C:\sigma\vdash S\coerce T\), then
\(\sigma\vdash_{\mathsf v} v:T\).
\end{lemma}

The proof is simultaneous induction on machine-term and value typing. For a
rule carrying a final coercion, \(C\) is composed after that coercion; the
\ruleref{term-value} case uses the value induction hypothesis.

\begin{lemma}[Typing interpretation]
\label{lem:semantic-typing}
If \(\mathcal R:\rho\realizes_\sigma\Gamma\) and
\(\Gamma\vdash t:T\), then
\[
  \sigma\vdash_{\mathsf{tm}}\rho(t):\rho(T).
\]
\end{lemma}

The proof is induction on the source typing derivation. The path case uses
typed path resolution. Abstraction and let suspend the premise beneath their binder
together with \(\mathcal R\). Pair introduction and application use their
syntax-directed result types with reflexive final coercions. In the subsumption case,
\cref{lem:coercion-construction} interprets the subtyping premise, and
\cref{lem:machine-postcomposition} composes the resulting coercion with the
machine-typing derivation.

\begin{lemma}[Body instantiation]
\label{lem:body-instantiation}
If \(B:\bodycl{\sigma}{z}{S}{t}{T}\) and
\(R_y:\sigma\realizes y:S\), then
\[
  \sigma\vdash_{\mathsf{tm}}\subst{t}{y}{z}:\subst{T}{y}{z}.
\]
\end{lemma}

Inverting \ruleref{body-source} yields a suspended source-typing derivation and a
closing substitution satisfying its outer context. The realization \(R_y\)
extends that substitution by \(z\mapsto y\); applying
\cref{lem:semantic-typing} to the suspended derivation proves the result.

For the empty context and store, \cref{lem:semantic-typing}, together with
\ruleref{kont-hole} and \ruleref{cfg-state}, gives
\(\vdash\mathsf{initial}(t):T\) from \(\ctxempty\vdash t:T\).

Typing interpretation is the bridge from source subtyping to the machine
invariant. Source subsumption becomes the final coercion carried by machine
typing. An elimination case applies that coercion before inverting the
realization of a stored value. Allocation constructs the introduction
realization for the fresh cell, applies the final coercion, and resumes the
continuation. Body and codomain premises are instantiated when beta reduction
or return supplies the bound location; a pair-member premise is instantiated
during coercion application after pair-realization inversion supplies the
stored first component.

\subsection{Progress}

Machine typing for a path \(p:T\), together with
\(\sigma\vdash p\resolve\loc x\), yields \(\sigma\realizes x:T\): construct
the singleton realization from the resolution and apply the final coercion.
This observation supplies canonical forms for applications. If the operator
has a function type, inversion of \ruleref{loc-fun} produces an abstraction at
the resolved location.

All remaining progress cases follow the outer term constructor. A
non-variable path resolves to its location; a variable is final under an empty
continuation or returns through a suspended let body. Values are final under
an empty continuation and allocate under a nonempty continuation. A let
expression pushes its body onto the continuation.

\begin{theorem}[Progress]
\label{thm:progress}
If \(\vdash c:T\), then either \(\final(c)\) or there is a state
\(c'\) such that \(c\sstep c'\).
\end{theorem}

\subsection{Preservation}

\begin{lemma}[Store extension]
\label{lem:store-extension}
If \(\sigma'\) is obtained by appending fresh bindings to \(\sigma\), then
resolution, store-induced path equality, signature conversion, and all
store-indexed semantic and machine-typing judgments over \(\sigma\) remain
valid over \(\sigma'\).
\end{lemma}

The named presentation chooses globally fresh locations, so weakening into
the extended store leaves existing names unchanged.

\begin{theorem}[One-step preservation]
\label{thm:preservation}
If \(\vdash c:T\) and \(c\sstep c'\), then \(\vdash c':T\).
\end{theorem}

The application case contains the essential dependent reasoning. Suppose
\(p\) resolves to \(x\) and \(q\) resolves to \(y\). Inverting
\ruleref{term-app} gives operator typing at \((z:S)\to U\), argument typing
at \(S\), and a final coercion
\(C_R:\sigma\vdash\subst{U}{q}{z}\coerce T\). Preservation constructs
typing for the reduct in the following order.
\begin{enumerate}
\item Apply the operator's final coercion to the singleton realization of
  \(x\). Inversion of \ruleref{loc-fun} then exposes the stored abstraction,
  a suspended body derivation \(B:\bodycl{\sigma}{z}{S_0}{t}{U_0}\), a domain
  coercion \(C_I:\sigma\vdash S\coerce S_0\), and an open codomain coercion
  \(F_O:\sigma;z:S\vdash U_0\coerce U\). Store lookup is functional, so
  this binding identifies the stored abstraction with the one used by the
  \ruleref{e-app} transition.
\item The argument typing and its resolution give
  \(R_y:\sigma\realizes y:S\). Coercion application gives
  \[
    \action{C_I}{\ruleuse{ref-loc}(R_y)}=\ruleuse{ref-loc}(R_y'),
    \qquad R_y':\sigma\realizes y:S_0.
  \]
\item Applying \cref{lem:body-instantiation} to \(B\) and the converted
  realization \(R_y'\) gives
  \[
    \sigma\vdash_{\mathsf{tm}}\subst{t}{y}{z}:\subst{U_0}{y}{z}.
  \]
\item Instantiating the open codomain coercion with the original argument
  realization gives
  \[
    C_F=\forcecod{F_O}{R_y},
    \qquad
    C_F:\sigma\vdash\subst{U_0}{y}{z}\coerce\subst{U}{y}{z}.
  \]
\item Since \(q\) and the location name \(y\) resolve to the same location,
  store-induced signature conversion gives a coercion \(C_q\) from
  \(\subst{U}{y}{z}\) to \(\subst{U}{q}{z}\). Composing the three
  coercions and applying \cref{lem:machine-postcomposition} gives
  \[
    \sigma\vdash_{\mathsf{tm}}\subst{t}{y}{z}:T.
  \]
\end{enumerate}
Thus the body is interpreted at the stored domain, while the open codomain
comparison is instantiated using the argument type ascribed by the function
realization.

For a path step, \(p\) resolves to \(x\). The target term has singleton type
\(\sing{x}\). Store-induced path equality gives a coercion from
\(\sing{x}\) to \(\sing{p}\), which composes with the path rule's final
coercion. Continuation typing for the frame pushed by a let step retains the
suspended body derivation.
At return, the returned location name's resolution and final coercion first
establish that its location realizes the frame's input type. Body
instantiation then types the suspended body at that location, and machine
subsumption composes the frame's final coercion. In the allocation case,
inversion exposes the value's introduction typing and final coercion.
After adding the fresh cell, the proof first weakens this value typing to the
extended store; the corresponding location-realization rule and coercion
application show that the new location realizes the value's ascribed type. It
then weakens the continuation-typing derivation, including its suspended body
derivation and final coercion, instantiates the body at the fresh location, and composes
the weakened frame coercion with the resulting machine typing.

\subsection{Finite executions}

Write \(\steps\) for the reflexive-transitive closure of \(\sstep\).
Iterating \cref{thm:preservation} carries configuration typing through
any finite execution. Applying \cref{thm:progress} at its last state gives the
standard non-stuckness statement.

\begin{theorem}[Finite preservation]
\label{thm:finite-preservation}
If \(\ctxempty\vdash t:T\) and
\(\mathsf{initial}(t)\steps c\), then \(\vdash c:T\).
\end{theorem}

\begin{theorem}[Type safety]
\label{thm:safety}
If \(\ctxempty\vdash t:T\) and
\(\mathsf{initial}(t)\steps c\), then \(c\) is final or can take a step.
\end{theorem}

The theorem quantifies over arbitrary finite prefixes. Divergent programs are
permitted; the last state of any finite prefix is not stuck.

\section{Mechanization and examples}
\label{sec:mechanization}

The Lean 4 development \citep{demoura2021lean4} uses intrinsically scoped
syntax. Types, terms, contexts, and stores are indexed by their available
variables or locations, and allocation extends the store scope. The named
presentation in this paper expresses the corresponding operations through
capture-avoiding substitution and globally fresh locations. Member signatures
are also indexed by their sort, so a term member cannot receive an interval
signature and a type member cannot receive a proper-type signature.

The formalization covers the complete calculus and abstract machine, typed
path resolution, subtyping interpretation, coercion application, machine
typing, progress, preservation, and \cref{thm:safety}. Allocation extends the
intrinsic store scope, so the preservation theorem returns the result type
weakened into that scope. The named presentation uses globally fresh
locations and leaves this renaming implicit. Closed regressions exercise
intersection introduction and both eliminations in one function application,
two incomparable same-label views of one stored record member, and the
derived merge of those views into one selectable member intersection. A
mixed two-cell regression uses one recursive plan to merge an older abstract
type member and, under the resulting record-prefix type, two different views
of a later term member, one of which depends on the prefix. Further closed
regressions merge two interval views of one stored abstract type member,
select the merged interval, and exercise both its lower and upper evidence;
the distinct-lower variant separately proves the joined-lower,
intersected-upper interval well formed and exercises union injection.

\subsection{Examples}

Intersections permit one function value to be used through two independently
derived views. Let
\[
  F=\top\to\top,
  \qquad S=\top\to F,
  \qquad L=F\to F,
  \qquad R=F.
\]
The value \(v=\lambda({\_}:\top).\lambda(y:\top).y\) has type \(S\).
Function subtyping gives \(S<:L\) by narrowing the outer domain and
\(S<:R\) by widening its codomain. Hence \ruleref{s-inter} gives
\(S<:\interty{L}{R}\), and the closed term
\[
  \mathbf{let}\ f=v\ \mathbf{in}\ f\,f
\]
has type \(F\): the operator occurrence uses \ruleref{s-inter-left} to view
\(f\) at \(L\), while the argument occurrence uses
\ruleref{s-inter-right} to view the same path at \(R\). At run time this is
ordinary function application; no intersection object is built.

The one-cell term-member plan makes the same two views useful behind one
record member. Write \(Q(X)=\pairty{x}{\top}{a}{X}\). Rule
\ruleref{m-pair}, with \ruleref{m-same} for the first component and
\ruleref{m-inter} for the member, followed by \ruleref{s-merge}, derives
\[
  \interty{Q(L)}{Q(R)}<:Q(\interty{L}{R}).
\]
The mechanized closed example first allocates \(v:S\) and the single pair
\(r=\pairval{v}{a}{v}:Q(S)\). Pair covariance lifts \(S<:L\) and
\(S<:R\) to the two record views, after which \ruleref{s-inter} assigns
\(r:\interty{Q(L)}{Q(R)}\). A path-only let gives one alias the merged type
\(Q(\interty{L}{R})\); ordinary \ruleref{p-sel} exposes the one stored
\(a\)-member at \(\interty{L}{R}\), and the same selected path supplies both
operator and argument of the application. The merge rule neither allocates a
record nor adds a projection from the outer intersection to path synthesis.

The mixed recursive regression demonstrates why the member premise of
\ruleref{m-pair} is formed under the merged first component. Define four
views of an older type-member cell:
\[
\begin{array}{ll}
 Q_0=\pairty{z}{\top}{A}{S..S},
 &Q_L=\pairty{z}{\top}{A}{S..L},\\[2pt]
 Q_R=\pairty{z}{\top}{A}{S..R},
 &Q_I=\pairty{z}{\top}{A}{S..\interty{L}{R}}.
\end{array}
\]
The corresponding outer term-member views are
\[
\begin{array}{rcl}
 P_0&=&\pairty{x}{Q_0}{f}{x.A},\\
 P_L&=&\pairty{x}{Q_L}{f}{x.A},\\
 P_R&=&\pairty{x}{Q_R}{f}{R},\\
 P_I&=&\pairty{x}{Q_I}{f}{\interty{x.A}{R}}.
\end{array}
\]
An inner pair plan preserves \(\top\), preserves the common interval lower
\(S\), and intersects the upper bounds, deriving
\(Q_L\mergeop Q_R\Longrightarrow Q_I\). The outer pair plan uses that whole
derivation for its first component. Its member subplan is checked under
\(x:Q_I\), where it intersects the genuinely dependent type \(x.A\) with
the direct view \(R\). Hence one recursive plan derives
\[
  P_L\mergeop P_R\Longrightarrow P_I,
  \qquad
  \interty{P_L}{P_R}<:P_I.
\]
This single step changes both axes of the outer pair; none of the former
one-axis principles expressed it.

The closed program allocates \(v:S\), then an older cell storing the exact
definition \(A=S\), and finally a newer cell storing \(f=v\) at the selected
type of the older cell. Pair and interval subtyping expose the resulting
physical spine at both \(P_L\) and \(P_R\), after which intersection
introduction and the recursive plan give one precise alias \(q:P_I\). Its
member has type
\[
  q.f:\interty{q.1.A}{R}.
\]
For the operator occurrence of \(q.f\), the left projection enters \(q.1.A\)
and its merged upper bound exits at \(L\). For the argument occurrence, the
right projection directly supplies \(R=F\). Thus the one stored path types
the self-application \(q.f\;q.f\), while the proof necessarily uses the
environment extension at the merged prefix described in
\cref{sec:coercion-action}.

An interval merge also places the intersection behind abstract type selection.
Write
\[
  P(U)=\pairty{x}{\top}{A}{S..U}.
\]
The lower plan preserves \(S\), the upper plan chooses an opaque intersection,
and \ruleref{m-pair} and \ruleref{s-merge} derive
\[
  \interty{P(L)}{P(R)}<:P(\interty{L}{R}).
\]
The mechanized closed example allocates the same value \(v:S\) and one pair
\(r=\pairval{v}{A}{S}\), whose type definition is exactly \(A=S\). Pair
covariance exposes \(r\) at \(P(L)\) and \(P(R)\), and intersection
introduction gives their conjunction. One path-only alias \(q\) receives the
merged type \(P(\interty{L}{R})\), so ordinary path synthesis exposes
\[
  q.A:S..\interty{L}{R}.
\]
In the body \(v\,v\), both occurrences of \(v\) first cross the selected
lower bound \(S<:q.A\). Selection of the merged upper bound followed by
\ruleref{s-inter-left} views the operator at \(L\); the corresponding
\ruleref{s-inter-right} path views the argument at \(R\). Thus the example
does not merely construct the merged interval: it uses the abstract selected
type to mediate the final application.

The union-lower plan is exercised with genuinely different lower bounds. Write
\[
  Q(X,Y)=\pairty{x}{\top}{A}{X..Y}.
\]
Rules \ruleref{j-union}, \ruleref{m-intv}, \ruleref{m-pair}, and
\ruleref{s-merge} derive
\[
  \interty{Q(\bot,L)}{Q(S,R)}
  <:
  Q(\unionty{\bot}{S},\interty{L}{R}).
\]
The mechanized closed program again stores the exact definition \(A=S\).
Interval and pair covariance expose that one pair at both \(Q(\bot,L)\) and
\(Q(S,R)\): the first view uses \(\bot<:S<:L\), while the second uses
\(S<:S<:R\). The merged target is well formed because
\[
  \unionty{\bot}{S}<:\interty{L}{R}.
\]
This derivation eliminates the union: the \(\bot\) arm is below both uppers by
\ruleref{s-bot}, and the \(S\) arm is below both by \(S<:L\) and \(S<:R\).
One precise alias \(q\) then exposes
\[
  q.A:\unionty{\bot}{S}..\interty{L}{R}.
\]
For each operand of \(v\,v\), \ruleref{s-union-right} first injects
\(v:S\) into \(\unionty{\bot}{S}\); the lower selection rule crosses into
\(q.A\), and the upper selection rule followed by the two intersection
projections supplies \(L\) for the operator and \(R\) for the argument. This
checks the full lower-join/upper-meet rule rather than obtaining it as an
unused type.

The following derivations illustrate \ruleref{s-pair} with a non-singleton
first component:
\[
  \pairty{x}{\top}{a}{\sing{x}}
  <:
  \pairty{x}{\top}{a}{\top}
\]
and
\[
  \pairty{x}{\top}{A}{\sing{x}..\sing{x}}
  <:
  \pairty{x}{\top}{A}{\bot..\top}.
\]
The first uses \ruleref{s-pair}, reflexivity for the first component, and
\ruleref{s-top} for the member. The second uses \ruleref{s-bounds} under the
pair binder; its lower, upper, and bound-consistency premises are respectively
\ruleref{s-bot}, \ruleref{s-top}, and reflexivity.

The interval derivation also occurs in the closed program
\[
\begin{aligned}
t_{\mathsf{int}}
  &={\mathbf{let}\ y=\lambda(x:\top).x\ \mathbf{in}\
      \pairval{y}{A}{\sing{y}}},
\\
\ctxempty&\vdash t_{\mathsf{int}}:
  \pairty{x}{\top}{A}{\bot..\top}.
\end{aligned}
\]
Pair introduction first assigns
\(\pairty{x}{\sing{y}}{A}{\sing{y}..\sing{y}}\). Under
\(x:\sing y\), \ruleref{s-widen} derives
\(\sing x<:\sing y\), \ruleref{s-alias} derives
\(\sing y<:\sing x\), and \ruleref{s-bounds} uses these premises with
reflexivity to derive
\(\sing y..\sing y<:\sing x..\sing x\). Hence one application of
\ruleref{s-pair} widens the first component to \(\top\) and expresses the
member interval as \(\sing{x}..\sing{x}\); a second application hides that
interval behind \(\bot..\top\). The mechanization specializes
\cref{thm:safety} to every finite execution from this term.

\section{Related work}
\label{sec:related}

\paragraph{Dependent object types.}
Foundational DOT systems account for abstract type members, dependent method
types, recursive objects, and rich subtyping
\citep{amin2014foundations,rompf2016dot}. Subsequent proofs factor the
metatheory in different ways, including inert or tight typing
\citep{rapoport2017simple} and definitional-interpreter techniques
\citep{amin2017definitional}. \lambdap{} removes recursive objects and studies
paths built from immutable dependent pairs. It retains controlled binary
intersections and unions with the three standard meet and join rules and a
plan-directed recursive merge judgment for aligned pair spines. The judgment
can simultaneously normalize a record prefix and a dependent later member,
but only when labels, member sorts, and field order align. Interval plans
preserve equal lowers or join distinct lowers and recursively merge the upper
bounds. Using a merged target in term typing separately requires that target
well formed. These connectives do not generally distribute through other type
constructors or participate directly in precise path synthesis, and the
proper-type union is not a sum type. The resulting runtime makes allocation
order available as a well-founded measure for the coercions generated by
dependent-pair subtyping.

pDOT supports types depending on paths of arbitrary length and uses singleton
types to track path equality \citep{rapoport2019pdot}. Its soundness proof
adapts syntactic DOT techniques and imposes conditions needed for path
lookup. \lambdap{} likewise separates singleton identity from abstract type
selection, while representing record traversal by pair lookup and proving
typed path resolution through store-indexed realization.

gDOT gives a step-indexed logical-relations model in Iris for a guarded DOT
calculus with recursive types and objects
\citep{giarrusso2020gdot,jung2018iris}. Its surface calculus exposes a later
type former and an explicit coercion term to control guarded recursive
reasoning. For the recursion-free \lambdap{} core, semantic coercions and the
allocation-order rank of referents suffice for well-foundedness.

\paragraph{Coercion interpretations.}
Coercion semantics translates subtyping or inheritance derivations into
explicit evidence. Transitivity becomes composition, and multiple derivations
of the same judgment raise a coherence question
\citep{breazutannen1991coercion,crary2000inclusive}. System FC makes type
equality evidence and casts explicit in a typed intermediate language and
proves erasure before execution \citep{sulzmann2007systemfc}. Our semantic
coercions follow a derivation-directed interpretation and preserve
realization at a fixed referent within a store. Type safety requires each
interpreted derivation to preserve realization; it does not require two
derivations of the same subtyping judgment to produce the same coercion. The
dependent cases suspend source derivations and interpret them only after a
concrete location is supplied.

Stratified semantic models have long used worlds or store levels to break
circular definitions involving higher-order state
\citep{ahmed2002stratified}. Location referents are ranked by their allocation
position in the immutable store, while type referents have rank zero. The
rank is used solely to justify recursive coercion application at referents
stored inside a pair.

\paragraph{Intervals, GADTs, and compilation to DOT.}
Type intervals also support a uniform account of higher-order bounded
abstraction and translucent definitions in
\(F^{\omega}_{..}\) \citep{stucki2021intervals}. \lambdap{} uses first-order
intervals as signatures of abstract type members represented by pairs and
retains explicit bound-consistency premises in selection, interval formation,
and interval subtyping.

Work on GADT reasoning in Scala and cDOT studies equality and subtyping
constraints recovered by pattern matching
\citep{parreaux2019gadt,boruchgruszecki2022cdot}. These systems refine the
typing assumptions within a branch after a constructor match, a facility that
would become relevant when extending \lambdap{} with term-level sums and pattern
matching. Type-preserving translations from class-based calculi to DOT study
how practical source-language constructs elaborate into a sound
path-dependent core \citep{martres2023compilation}. Such a translation would
also provide a setting in which to assess the coverage of the dependent-pair
encodings used here.

\section{Conclusion}

We have proved type safety for \lambdap{} with arbitrary stable paths,
singleton types, controlled binary intersections and unions, abstract type
members with interval signatures, explicit subtyping transitivity, and
covariance of both components of dependent pairs.
The semantic interpretation records which signatures a location or stored
type definition realizes. Interpreting a subtyping derivation produces a
coercion that preserves the referent while changing its realization.

The dependent-pair covariance case retains the member-subtyping premise until inversion
reveals the stored first component. Instantiating the premise at that location
yields a member coercion, which is applied to the pair's member referent. This
coercion may be larger than the enclosing pair coercion; recursive application
remains well founded because both component referents have lower allocation
rank than the pair cell. Together with typed path resolution and the
syntax-directed machine invariant, the coercion interpretation establishes
progress, preservation, and type safety for every finite execution prefix.
The meet/join extension is intentionally limited to the three standard rules
for each connective and one plan-directed aligned merge judgment. A plan may
retain an opaque intersection or recursively merge same-label pair spines;
the member recursion is interpreted under the merged first-component binder.
Interval plans preserve a shared lower or form its union and recursively merge
their uppers. The four former one-cell principles are derived instances of
this judgment. A plan does not discharge target well-formedness, so term
subsumption still requires a separate proof, including interval consistency.
The calculus has no path-typing projection through either connective, merge
for unequal labels or differently ordered fields, canonical merge normal
form, other distributivity, dedicated term introduction or elimination form,
sum values, or operational merge rule; extending any of those dimensions
would be a separate calculus change.

\bibliographystyle{ACM-Reference-Format}
\bibliography{../../references}

\end{document}
